\chaptermark{Exceptional publication of the continuum}
\label{chap:83-20}

\input{ampliaciones_sucesoras_20260824/parte_i_ii/tex/bloques/salto_espectral_4_2_6_54}

% Unidad U018. Paley, extension cuadratica finita y codigo ternario.
% Redaccion autonoma desde la autoridad material 1063.

\section{Paley conference form, finite quadratic extension and self-dual ternary code}
\label{p12-83-c20-s1:chap:paley-extension-cuadratica-codigo-ternario}

The tetrad \(B_K\) constructed in the preceding unit does not yet have a
Witt star: to define it, one must first construct the incidence system that
contains its hexads. This unit performs that step from the
projective gauge of the six units modulo nine, obtains an internal square
root of \(-I_6\), and constructs the ternary graph code of
length twelve. No identification follows from a mere equality of
cardinalities.

\subsection{Projective gauge of the \(6\) units modulo \(9\)}

Let
\[
 \mathcal U_9=(\mathbb Z/9\mathbb Z)^\times
 =\{1,2,4,8,7,5\},
\]
ordered by the successive powers of \(2\) modulo \(9\). Fijamos the
coordinate bijection
\begin{equation}
 \theta:\mathcal U_9\longrightarrow\mathbb P^1(\mathbb F_5),
 \qquad
 \begin{array}{c|cccccc}
 u&1&2&4&8&7&5\\ \hline
 \theta(u)&\infty&0&1&2&3&4
 \end{array}.
 \label{p12-83-c20-s1:eq:calibre-unidades-paley}
\end{equation}
The map \(\theta\) is a coordinate gauge. It is not asserted to be a homomorphism between multiplication in \(\mathcal U_9\) and an operation on \(\mathbb P^1(\mathbb F_5)\).

Let
\[
 \chi:\mathbb F_5\longrightarrow\{-1,0,1\}
\]
the quadratic character:
\[
 \chi(0)=0,
 \qquad
 \chi(1)=\chi(4)=1,
 \qquad
 \chi(2)=\chi(3)=-1.
\]
In the order \((\infty,0,1,2,3,4)\), we define
\begin{align}
 (C_P)_{\infty,x}&=(C_P)_{x,\infty}=1,
 &&x\in\mathbb F_5,\label{p12-83-c20-s1:eq:paley-fila-infinito}\\
 (C_P)_{x,y}&=\chi(x-y),
 &&x,y\in\mathbb F_5, x\neq y,\label{p12-83-c20-s1:eq:paley-entradas-finitas}\\
 (C_P)_{x,x}&=0.
 \label{p12-83-c20-s1:eq:paley-diagonal}
\end{align}
Thus,
\begin{equation}
 C_P=
 \begin{pmatrix}
 0& 1& 1& 1& 1& 1\\
 1& 0& 1&-1&-1& 1\\
 1& 1& 0& 1&-1&-1\\
 1&-1& 1& 0& 1&-1\\
 1&-1&-1& 1& 0& 1\\
 1& 1&-1&-1& 1& 0
 \end{pmatrix}.
 \label{p12-83-c20-s1:eq:matriz-conferencia-paley-autonoma}
\end{equation}

\begin{lemma}[Quadratic correlation sum]
\label{p12-83-c20-s1:lem:correlacion-cuadratica-f5}
For \(a\neq b\) in \(\mathbb F_5\),
\begin{equation}
 \sum_{t\in\mathbb F_5}
 \chi(t-a)\chi(t-b)=-1.
 \label{p12-83-c20-s1:eq:correlacion-cuadratica-f5}
\end{equation}
\end{lemma}

\begin{proof}
The substitution \(u=(t-a)/(b-a)\) is to bijection of \(\mathbb F_5\).
Ifnce \(\chi((b-a)^2)=1\), the sum  reduce to
\[
 \sum_{u\in\mathbb F_5}\chi(u(u-1)).
\]
Evaluation of \(u=0,1,2,3,4\) gives, respectively,
\(0,0,1,-1,-1\), whose sum is \(-1\).
\end{proof}

\begin{proposition}[Conference Identity]
\label{p12-83-c20-s1:prop:paley-identidad-conferencia-autonoma}
The matrix \(C_P\) is symmetric and satisfies
\begin{equation}
 \boxed{C_PC_P^{\mathsf T}=5I_6.}
 \label{p12-83-c20-s1:eq:paley-identidad-conferencia-autonoma}
\end{equation}
\end{proposition}

\begin{proof}
Each row contains five entries of modulus one and one zero entry; its
squared norm is five. For two distinct finite rows, the contribution
of the coordinate \(\infty\) is \(+1\), whereas
\cref{p12-83-c20-s1:lem:correlacion-cuadratica-f5} gives \(-1\) for the five finite
coordinates. Their inner product is zero. The sum of a nontrivial character
over \(\mathbb F_5\) is zero, so the row corresponding to
\(\infty\) is also orthogonal to every finite row.
\end{proof}

\subsection{ternary reduction and internal quadratic root of \texorpdfstring{\(-I_6\)}{-I₆}}

We reduce \(C_P\) modulo tris, with \(-1\equiv2\pmod3\):
\begin{equation}
 A_W=
 \begin{pmatrix}
 0&1&1&1&1&1\\
 1&0&1&2&2&1\\
 1&1&0&1&2&2\\
 1&2&1&0&1&2\\
 1&2&2&1&0&1\\
 1&1&2&2&1&0
 \end{pmatrix}
 \in M_6(\mathbb F_3).
 \label{p12-83-c20-s1:eq:matriz-paley-witt-ternaria}
\end{equation}

\begin{theorem}[Internal square root of \(-I_6\)]
\label{p12-83-c20-s1:thm:aw-cuadrado-menos-identidad-autonomo}
We have
\begin{equation}
 \boxed{A_W^2=-I_6.}
 \label{p12-83-c20-s1:eq:aw-cuadrado-menos-identidad-autonomo}
\end{equation}
In particular,
\[
 A_W^{-1}=-A_W.
\]
\end{theorem}

\begin{proof}
The \cref{p12-83-c20-s1:prop:paley-identidad-conferencia-autonoma} gives
\[
 A_WA_W^{\mathsf T}
 \equiv5I_6
 \equiv2I_6
 =-I_6
 \pmod3.
\]
Since \(A_W=A_W^{\mathsf T}\), \(A_W^2=-I_6\) is obtained.
\end{proof}

The preceding identity is the algebraic residence of the root of \(-1\)
in this corridor. No proceeds of the factor \(1/4\) of the Hadamard inversion:
the former is a rational scalar forced by \(H_4^2=4I_4\); this is one
quadratic structure on characteristic three.

\subsection{Extension quadratic to \texorpdfstring{\(\mathbb F_9\)}{F₉}}

We define
\begin{equation}
 \mathbb F_9=\mathbb F_3[\iota]/(\iota^2+1),
 \qquad \iota^2=-1.
 \label{p12-83-c20-s1:eq:definicion-f9-autonoma}
\end{equation}
The polynomial \(X^2+1\) has no roots in \(\mathbb F_3\), so that the
quotient is a field of nine elements. For
\(V=\mathbb F_3^6\), let
\[
 V_9=V\otimes_{\mathbb F_3}\mathbb F_9.
\]
The operators
\begin{equation}
 P_\iota=\frac12(I-\iota A_W),
 \qquad
 P_{-\iota}=\frac12(I+\iota A_W)
 \label{p12-83-c20-s1:eq:proyectores-f9-paley}
\end{equation}
are well defined because \(2\) is invertible in \(\mathbb F_3\).

\begin{proposition}[Conjugate projectors]
\label{p12-83-c20-s1:prop:proyectores-conjugados-paley}
The following hold
\[
 P_\iota^2=P_\iota,
 \qquad
 P_{-\iota}^2=P_{-\iota},
 \qquad
 P_\iota P_{-\iota}=0,
 \qquad
 P_\iota+P_{-\iota}=I.
\]
\end{proposition}

\begin{proof}
Usando \(\iota^2=-1\) and \(A_W^2=-I\),
\[
 (\iota A_W)^2=I.
\]
The four identities are then the elementary identities for the projectors
\((I\mp\iota A_W)/2\) of an involution.
\end{proof}

\begin{theorem}[Conjugate spectral decomposition]
\label{p12-83-c20-s1:thm:descomposicion-espectral-f9-paley}
Let
\[
 E_{\pm\iota}=\ker(A_W\mp\iota I).
\]
Then
\begin{equation}
 V_9=E_\iota\oplus E_{-\iota},
 \qquad
 \dim_{\mathbb F_9}E_\iota
 =\dim_{\mathbb F_9}E_{-\iota}=3.
 \label{p12-83-c20-s1:eq:descomposicion-espectral-f9-paley}
\end{equation}
\end{theorem}

\begin{proof}
The minimal polynomial of \(A_W\) divides \(X^2+1\), which splits over
\(\mathbb F_9\) with simple roots \(\pm\iota\). The projectors of
\cref{p12-83-c20-s1:prop:proyectores-conjugados-paley} produce the direct sum. Since
\(A_W\) has coefficients in \(\mathbb F_3\), the Frobenius automorphism
exchanges the two eigenspaces; their dimensions are equal and
sum to six.
\end{proof}

This construction dota \(\mathbb F_3^6\) of istructura of modulo of rango
tris over \(\mathbb F_9\), by means of
\begin{equation}
 (a+b\iota)\cdot v:=av+b(vA_W).
 \label{p12-83-c20-s1:eq:accion-f9-sobre-f3-seis}
\end{equation}
The equality \(3^6=9^3\) is thus accompanied by an explicit action; it is not
used as a cardinal argument.

\subsection{Ternary graph code of length \(12\)}

We define
\begin{equation}
 c:\mathbb F_3^6\longrightarrow\mathbb F_3^{12},
 \qquad
 c(q)=(q,qA_W),
 \label{p12-83-c20-s1:eq:aplicacion-codigo-grafico-witt}
\end{equation}
and
\begin{equation}
 C_W:=\operatorname{im}c.
 \label{p12-83-c20-s1:eq:definicion-codigo-witt-ternario}
\end{equation}
Its generating matrix is
\[
 G_W=(I_6\mid A_W).
\]

\begin{theorem}[Injectivity and self-duality]
\label{p12-83-c20-s1:thm:cw-inyectivo-autodual}
The map \(c\) is injective, \(\dim C_W=6\) and
\begin{equation}
 \boxed{C_W=C_W^\perp.}
 \label{p12-83-c20-s1:eq:cw-autodual}
\end{equation}
\end{theorem}

\begin{proof}
The first half of \(c(q)\) is \(q\), so
\(\operatorname{pr}_1\circ c=\operatorname{id}\); this proves
injectivity and the dimension. Moreover,
\[
 G_WG_W^{\mathsf T}
 =I_6+A_WA_W^{\mathsf T}
 =I_6+2I_6=0
\]
in \(\mathbb F_3\). Therefore, \(C_W\subseteq C_W^\perp\). Both spaces
have dimension six, so they are equal.
\end{proof}

For \(r=0,\ldots,12\), let
\[
 N_r:=\#\{q\in\mathbb F_3^6:
            \operatorname{wt}(q,qA_W)=r\}.
\]
The enumeration of the \(3^6=729\) entries produces
\begin{equation}
 N_0=1,
 \qquad
 N_6=264,
 \qquad
 N_9=440,
 \qquad
 N_{12}=24,
 \label{p12-83-c20-s1:eq:censo-pesos-cw}
\end{equation}
and \(N_r=0\) in the remaining weights.

\begin{theorem}[Enumerator and minimum distance]
\label{p12-83-c20-s1:thm:enumerador-cw-autonomo}
The code \(C_W\) has parameters
\[
 \boxed{C_W=[12,6,6]_3}
\]
and weight enumerator
\begin{equation}
 \boxed{W_{C_W}(y)=1+264y^6+440y^9+24y^{12}.}
 \label{p12-83-c20-s1:eq:enumerador-pesos-cw}
\end{equation}
\end{theorem}

\begin{proof}
The dimension and self-duality have already been proved. The census
\eqref{p12-83-c20-s1:eq:censo-pesos-cw} traverses the \(729\) words without sampling. Its
least nonzero weight is six, which is the minimum distance of a linear code.
The sum \(1+264+440+24=729\) checks that no word was omitted.
\end{proof}

\subsection{Nature of the structure obtained}

Three distinct objects and three typed arrows have been constructed:
\[
 \mathcal U_9
 \xrightarrow{\theta}\mathbb P^1(\mathbb F_5)
 \xrightarrow{\chi}C_P
 \xrightarrow{\bmod3}A_W,
\]
\[
 A_W^2=-I_6
 \quad\Longrightarrow\quad
 \mathbb F_3^6\cong\mathbb F_9^3
 \quad\Longrightarrow\quad
 C_W=\{(q,qA_W):q\in\mathbb F_3^6\}.
\]
The first arrow is a gauge, the second uses the quadratic character, and the third is coefficient reduction. The structure over
\(\mathbb F_9\) is finite; it is not identified with a real or analytic complex structure.

\subsection{Obstructions in the Paley–\(C_W\) corridor: bijection, \(A_W^2=-I_6\), self-duality, and census of \(729\) words}

\begin{criterion}[Falsification of the Paley--\(C_W\) corridor]
The construction is refuted if any of the following occurs:
\begin{enumerate}
 \item \(\theta\) is not bijective or is used as a homomorphism without proof;
 \item the sum \eqref{p12-83-c20-s1:eq:correlacion-cuadratica-f5} differs from \(-1\);
 \item \(C_PC_P^{\mathsf T}\neq5I_6\);
 \item \(A_W^2\neq-I_6\);
 \item the projectors \eqref{p12-83-c20-s1:eq:proyectores-f9-paley} are not
       complementary;
 \item the map \(q\mapsto(q,qA_W)\) loses injectivity or
       self-duality;
 \item the census does not sum to \(729\), differs from
       \eqref{p12-83-c20-s1:eq:censo-pesos-cw}, or presents a word of weight between one
       and five;
 \item \(\mathbb F_3^6\cong\mathbb F_9^3\) is inferred only from
       \(3^6=9^3\), omitting the action \eqref{p12-83-c20-s1:eq:accion-f9-sobre-f3-seis}.
\end{enumerate}
\end{criterion}


% Unidad U019. Sistema de Witt, campo de estrellas y grupo M12.
% Redaccion autonoma desde la autoridad material 1063.

\section{Ternary Steiner system, Witt-star involutions and generation of \texorpdfstring{\(M_{12}\)}{M12}}
\label{p12-83-c20-s2:chap:sistema-witt-estrellas-m12}

The preceding unit constructed the self-dual code
\(C_W=[12,6,6]_3\). Its words of weight six are now projectivized,
the Steiner property is proved, and only then is the star
associated with the tetrad \(B_K\) defined. This order prevents an incidence
figure from being turned into an object added by visual similarity.

\subsection{Supports projective of weight \(6\)}

Let
\begin{equation}
 \mathcal H_W
 :=\{\operatorname{supp}(c):c\in C_W,
                              \operatorname{wt}(c)=6\}.
 \label{p12-83-c20-s2:eq:familia-soportes-peso-seis}
\end{equation}
Each support occurs for the word pair \(\{c,-c\}\). Indeed, if two
weight-six words with the same support were not proportional, one could
scale one of them to cancel a coordinate upon subtraction; this would yield
a nonzero word of \(C_W\) of weight less than six, contrary to the minimum
distance. Since neither is there a nonzero word equal to its opposite in
characteristic three, every projective orbit has exactly two elements.
Therefore,
\begin{equation}
 |\mathcal H_W|=\frac{264}{2}=132.
 \label{p12-83-c20-s2:eq:numero-hexadas-cw}
\end{equation}

For \(T\in\binom{[12]}5\), define
\begin{equation}
 \nu(T):=\#\{H\in\mathcal H_W:T\subseteq H\}.
 \label{p12-83-c20-s2:eq:incidencia-cinco-hexada}
\end{equation}
The exact enumeration of the
\[
 \binom{12}{5}=792
\]
subsets of five points produce
\begin{equation}
 \nu(T)=1
 \qquad
 \text{for every }T\in\binom{[12]}5.
 \label{p12-83-c20-s2:eq:unicidad-cinco-en-hexada}
\end{equation}

\begin{theorem}[Ternary Witt System]
\label{p12-83-c20-s2:thm:w12-autonomo}
The incidence system
\[
 W_{12}:=([12],\mathcal H_W)
\]
is the Steiner system
\begin{equation}
 \boxed{W_{12}=S(5,6,12).}
 \label{p12-83-c20-s2:eq:w12-steiner-autonomo}
\end{equation}
\end{theorem}

\begin{proof}
The equality \eqref{p12-83-c20-s2:eq:unicidad-cinco-en-hexada} is precisely the
Steiner property. The number of blocks is also recovered from the double
count of pairs \((T,H)\) with \(T\subset H\):
\[
 |\mathcal H_W|\binom65=\binom{12}{5},
 \qquad
 |\mathcal H_W|=\frac{792}{6}=132.
\]
\end{proof}

\subsection{Residue of a tetrad and decomposition into \(4\) incident components}

Let \(B\in\binom{[12]}4\). In an \(S(5,6,12)\) system, the number of
hexads containing \(B\) is
\begin{equation}
 \lambda_4
 =\frac{\binom{12-4}{5-4}}{\binom{6-4}{5-4}}
 =\frac82=4.
 \label{p12-83-c20-s2:eq:lambda-cuatro-w12}
\end{equation}
We write these four hexads as
\[
 H_i=B\sqcup P_i,
 \qquad |P_i|=2,
 \qquad i=1,\ldots,4.
\]

\begin{lemma}[Residual partition of a tetrad]
\label{p12-83-c20-s2:lem:particion-residual-tetrada-witt}
The four pairs \(P_1,\ldots,P_4\) are disjoint and
\begin{equation}
 [12]\setminus B
 =P_1\sqcup P_2\sqcup P_3\sqcup P_4.
 \label{p12-83-c20-s2:eq:particion-residual-tetrada-witt}
\end{equation}
\end{lemma}

\begin{proof}
If two residual pairs shared a point \(p\), the corresponding
hexads would contain the same subset \(B\cup\{p\}\) of five points,
in contradiction with the Steiner uniqueness. The four pairs are, therefore,
disjoint; contain eight points and agotan the complement of \(B\).
\end{proof}

\begin{definition}[Star and Witt star involution]
\label{p12-83-c20-s2:def:estrella-witt-involucion}
The star centered at \(B\) is
\begin{equation}
 \mathscr S_B
 :=\{H\in\mathcal H_W:B\subseteq H\}
 =\{B\sqcup P_1,\ldots,B\sqcup P_4\}.
 \label{p12-83-c20-s2:eq:estrella-witt-general}
\end{equation}
If \(P_i=\{p_i^-,p_i^+\}\), its involution is
\begin{equation}
 \sigma_B:=\prod_{i=1}^{4}(p_i^-\ p_i^+).
 \label{p12-83-c20-s2:eq:involucion-estrella-witt-general}
\end{equation}
The permutation \(\sigma_B\) fixes \(B\) pointwise and exchanges the
ends of each residual petal.
\end{definition}

The mathematical term used henceforth is \emph{Witt star} or
\emph{Witt-star involution}. It is not a Euclidean
star.

\subsection{Star selected by the tetrad of the seal}

The dodecaphasic unit produced, without using Steiner's design,
\[
 B_K=\{4,6,9,10\}.
\]
Now that \(W_{12}\) has been constructed, its star is well defined. The
four incident hexads are
\begin{align}
 H_{13}&=B_K\sqcup\{1,3\}
       =\{1,3,4,6,9,10\},
 \label{p12-83-c20-s2:eq:hexada-petalo-13}\\
 H_{2,11}&=B_K\sqcup\{2,11\}
       =\{2,4,6,9,10,11\},
 \label{p12-83-c20-s2:eq:hexada-petalo-211}\\
 H_{5,7}&=B_K\sqcup\{5,7\}
       =\{4,5,6,7,9,10\}
       =H_\alpha^-,
 \label{p12-83-c20-s2:eq:hexada-petalo-57}\\
 H_{8,12}&=B_K\sqcup\{8,12\}
       =\{4,6,8,9,10,12\}.
 \label{p12-83-c20-s2:eq:hexada-petalo-812}
\end{align}
Consequently,
\begin{equation}
 \boxed{
 \sigma_{B_K}=(1\ 3)(2\ 11)(5\ 7)(8\ 12).}
 \label{p12-83-c20-s2:eq:involucion-estrella-witt-bk}
\end{equation}

\begin{figure}[htbp]
\centering
\begin{tikzpicture}[
  >=Stealth,
  font=\small,
  core/.style={draw,double,rounded corners=2mm,align=center,
    minimum width=4.0cm,minimum height=1.3cm,fill=black!4},
  petal/.style={draw,rounded corners=2mm,align=center,
    minimum width=3.55cm,minimum height=1.08cm,fill=black!2},
  negative/.style={petal,line width=1.1pt,double,fill=black!10},
  arr/.style={->,thick}
]
\node[core] (B) at (0,0)
 {\(\displaystyle B_K=\{4,6,9,10\}\)\\
  points fixed by \(\sigma_{B_K}\)};
\node[petal] (P13) at (0,3.0)
 {\(P_1=\{1,3\}\), \((1\ 3)\)\\
  \(H_{13}=B_K\sqcup P_1\)};
\node[petal] (P211) at (5.25,0)
 {\(P_2=\{2,11\}\), \((2\ 11)\)\\
  \(H_{2,11}=B_K\sqcup P_2\)};
\node[negative] (P57) at (0,-3.0)
 {\(P_3=\{5,7\}\), \((5\ 7)\)\\
  \(H_{5,7}=H_\alpha^-\)};
\node[petal] (P812) at (-5.25,0)
 {\(P_4=\{8,12\}\), \((8\ 12)\)\\
  \(H_{8,12}=B_K\sqcup P_4\)};
\draw[arr] (B)--(P13);
\draw[arr] (B)--(P211);
\draw[arr] (B)--(P57);
\draw[arr] (B)--(P812);
\node[align=center,font=\footnotesize] at (0,-4.35)
 {\(\mathscr S_{B_K}
   =\{H_{13},H_{2,11},H_{5,7},H_{8,12}\}\).\\
  The diagram represents incidence, not geometric distance.};
\end{tikzpicture}
\caption{Witt star centered at \(B_K\). The center is the fixed tetrad;
the four petals are the residual pairs that determine the
involution.}
\label{p12-83-c20-s2:fig:estrella-witt-bk-autonoma}
\end{figure}

\subsection{Global field of the 495 stars}

There exists a star for each tetrad, so
\begin{equation}
 \#\{\mathscr S_B:B\in\tbinom{[12]}4\}
 =\binom{12}{4}=495.
 \label{p12-83-c20-s2:eq:numero-estrellas-witt}
\end{equation}
The exact census verifies
\begin{equation}
 \sigma_B(\mathcal H_W)=\mathcal H_W
 \qquad
 \text{for every }B\in\binom{[12]}4.
 \label{p12-83-c20-s2:eq:estrellas-conservan-hexadas}
\end{equation}
We define
\begin{equation}
 \Gamma_\star
 :=\langle\sigma_B:B\in\tbinom{[12]}4\rangle
 \leq\operatorname{Aut}(W_{12}).
 \label{p12-83-c20-s2:eq:grupo-generado-estrellas-witt}
\end{equation}

A generating set of six tetrads is
\[
 1234,
 \quad1235,
 \quad1236,
 \quad1245,
 \quad1345,
 \quad2345.
\]
The orders of the subgroups obtained by adding their involutions successively are
\begin{equation}
 1, 2, 6, 18, 360, 7920, 95040.
 \label{p12-83-c20-s2:eq:cadena-ordenes-generadores-m12}
\end{equation}

\begin{theorem}[Generation of the Mathieu group \(M_{12}\)]
\label{p12-83-c20-s2:thm:495-estrellas-m12-autonomo}
We have
\begin{equation}
 \boxed{
 \langle\sigma_B:B\in\tbinom{[12]}4\rangle
 =\operatorname{Aut}(W_{12})
 \cong M_{12}.}
 \label{p12-83-c20-s2:eq:estrellas-generan-m12}
\end{equation}
\end{theorem}

\begin{proof}
Each \(\sigma_B\) preserves the \(132\) hexads, hence
\(\Gamma_\star\leq\operatorname{Aut}(W_{12})\). The Schreier enumeration of
the six indicated generators yields order \(95040\). The automorphism group
of \(S(5,6,12)\) is \(M_{12}\) and has that same order; therefore
the inclusion is an equality.
\end{proof}

An isolated star generates only
\[
 \langle\sigma_B\rangle\cong C_2.
\]
The conclusion about \(M_{12}\) belongs to the global field of the
\(495\) involutions, not to \(\sigma_{B_K}\) considered in isolation.

\subsection{Action on the augmentation module}

Let
\[
 V_{11}=\mathbf1^\perp\subset\mathbb R^{12}.
\]
The involution \(\sigma_B\) fixes the four coordinates of \(B\) and acts by four transpositions on its complement. In \(\mathbb R^{12}\) it has eight directions with eigenvalue \(+1\) and four with eigenvalue \(-1\). The constant direction belongs to the first eigenspace. Therefore,
\begin{equation}
 \operatorname{Spec}(\sigma_B|_{V_{11}})
 =\{1^{[7]},(-1)^{[4]}\},
 \qquad
 \frac1{11}\operatorname{Tr}(\sigma_B|_{V_{11}})=\frac3{11}.
 \label{p12-83-c20-s2:eq:espectro-involucion-estrella-aumentacion}
\end{equation}
It is a signed involution; it is not a positive projector of rank three.

\subsection{Obstructions concerning \(W_{12}\), the field of \(495\) star involutions and the generation of \(M_{12}\)}

\begin{criterion}[Refutation of \(W_{12}\), its stars, and \(M_{12}\)]
The construction is refuted if any of the following facts occurs:
\begin{enumerate}
 \item the projectivization of the \(264\) words of weight six does not produce
       \(132\) supports;
 \item there exists a subset of five points contained in zero or in more of
       one hexad;
 \item a tetrad does not belong to exactly four hexads;
 \item the four residual pairs do not partition the complement of the
       tetrad;
 \item the petals of \(B_K\) are not
       \(\{1,3\},\{2,11\},\{5,7\},\{8,12\}\);
 \item any of the \(495\) involutions fails to preserve
       \(\mathcal H_W\);
 \item the generated group does not have order \(95040\);
 \item \(M_{12}\) is attributed to a single star instead of the global field of involutions;
 \item  introduces the istrella before of constructing \(W_{12}\).
\end{enumerate}
\end{criterion}


% Unidad U020. Prolongacion binaria W12--W24 y M24.
% Redaccion autonoma desde la autoridad material 1063.

\section{Binary extension relative to a dodecad, system \texorpdfstring{\(W_{24}\)}{W24} and group \texorpdfstring{\(M_{24}\)}{M24}}
\label{p12-83-c20-s3:chap:dodecada-w24-m24}

The system \(W_{12}\) constructed over twelve phases admits a binary prolongation
when an extended Golay code, a dodecad, and an incidence
isomorphism are declared. These data are relative inputs of this branch:
they are not deduced solely from \(C_W\), and they do not intervene in the ternary lattice
branch that will lead to the Leech lattice.

\subsection{Golay binary code and octads}

Let
\[
 \mathcal G_{24}\subset\mathbb F_2^{24}
\]
the extended binary Golay code. Its weights are
\begin{equation}
 0, 8, 12, 16, 24.
 \label{p12-83-c20-s3:eq:pesos-golay-binario-extendido}
\end{equation}
The supports of the weight eight words form the Steiner system
\begin{equation}
 W_{24}=S(5,8,24).
 \label{p12-83-c20-s3:eq:w24-steiner-octadas}
\end{equation}
Its blocks are called \emph{octads}. Fix a dodecad \(D\), that is,
the support of a word of weight twelve, and let \(\mathcal O_{24}\)
be the family of octads.

We define
\begin{equation}
 \mathcal H(D)
 :=\{O\cap D:O\in\mathcal O_{24},\ |O\cap D|=6\}.
 \label{p12-83-c20-s3:eq:hexadas-residuales-dodecada}
\end{equation}

\begin{theorem}[Residual design of a dodecada]
\label{p12-83-c20-s3:thm:diseno-residual-dodecada-autonomo}
The residual system is
\begin{equation}
 \boxed{(D,\mathcal H(D))\cong S(5,6,12).}
 \label{p12-83-c20-s3:eq:diseno-residual-dodecada-w12}
\end{equation}
Moreover, each \(H\in\mathcal H(D)\) has one unique octad \(O_H\) such that
\begin{equation}
 O_H\cap D=H.
 \label{p12-83-c20-s3:eq:elevacion-unica-hexada-octada}
\end{equation}
\end{theorem}

\begin{proof}
Let \(A\subset D\), with \(|A|=5\). In \(W_{24}=S(5,8,24)\) there is a
unique octad \(O_A\) containing \(A\). If \(j=|O_A\cap D|\), the binary
sum of the characteristic words of \(O_A\) and \(D\) belongs to
\(\mathcal G_{24}\) and has weight
\[
 8+12-2j=20-2j.
\]
Since \(A\subseteq O_A\cap D\), one has \(5\leq j\leq8\). The four
possible weights \(10,8,6,4\), combined with
\eqref{p12-83-c20-s3:eq:pesos-golay-binario-extendido}, force \(j=6\). Thus every
five-point subset of \(D\) belongs to a unique residual hexad.

If two octads had the same hexadic intersection with \(D\), both
would contain the same five points of that hexad, contradicting
uniqueness in \(S(5,8,24)\). This also proves uniqueness of the
lift.
\end{proof}

\subsection{Relative scale between the HMT design and the dodecagon}

To relate the preceding construction to the design obtained from
\(C_W\), one must explicitly declare
\begin{equation}
 (D,\ell),
 \qquad
 \ell:W_{12}^{\rm HMT}
 \xrightarrow{\sim}(D,\mathcal H(D)),
 \label{p12-83-c20-s3:eq:dato-relativo-dodecada-alineamiento}
\end{equation}
where \(\ell\) is an incidence isomorphism. If \(\ell_*\) denotes its
action on blocks, we define
\begin{equation}
 \widehat\iota_{D,\ell}
 :=\iota_D\circ\ell_*:
 W_{12}^{\rm HMT}\longrightarrow W_{24},
 \qquad
 \widehat\iota_{D,\ell}(H)=O_{\ell(H)}.
 \label{p12-83-c20-s3:eq:elevacion-tipificada-w12-w24}
\end{equation}
The composition \(\widehat\iota_{D,\ell}\) avoids applying
\(\iota_D\), whose domain is \(\mathcal H(D)\), directly to a hexad
that still lies in \([12]_{\rm HMT}\).

\begin{proposition}[Injectivity of the Relative Lift]
\label{p12-83-c20-s3:prop:inyectividad-elevacion-w12-w24}
The map \(\widehat\iota_{D,\ell}\) is injective and preserves the
incidence between points and hexads after applying \(\ell\).
\end{proposition}

\begin{proof}
The isomorphism \(\ell\) is bijective on blocks. By the
\cref{p12-83-c20-s3:thm:diseno-residual-dodecada-autonomo}, each residual hexad has a
unique lifting octad. Two HMT hexads with the same image would have the same
residual hexad and, by bijectivity of \(\ell\), would be equal.
\end{proof}

There is no absolute choice of \((D,\ell)\) prior to this gauge. The
invariant object is the class of the construction under simultaneous change of
dodecad and alignment.

\subsection{Incidence law \texorpdfstring{\(4+1\)}{4+1} on a tetrad}

For a tetrad \(B\subset D\), the parameters derived from the two
Steiner systems are
\begin{equation}
 \lambda_4(W_{12})
 =\frac{\binom81}{\binom21}=4,
 \qquad
 \lambda_4(W_{24})
 =\frac{\binom{20}1}{\binom41}=5.
 \label{p12-83-c20-s3:eq:ley-cuatro-mas-uno-witt}
\end{equation}
Therefore, the five octads containing \(B\) decompose uniquely
as
\begin{equation}
 \{O_H:H\in\mathcal H(D),\ B\subset H\}
 \sqcup\{O_B^\perp\}.
 \label{p12-83-c20-s3:eq:descomposicion-cinco-octadas}
\end{equation}
The four first satisfy \(|O_H\cap D|=6\); the quinta satisfies
\begin{equation}
 O_B^\perp\cap D=B.
 \label{p12-83-c20-s3:eq:octada-transversal-tetrada}
\end{equation}

For \(B=B_K\), the four residual lifts correspond to the
petals
\[
 \{1,3\},
 \qquad
 \{2,11\},
 \qquad
 \{5,7\},
 \qquad
 \{8,12\},
\]
and the fifth octad is the transversal \(O_{B_K}^\perp\). This law explains
the transition from four hexads to five octads without identifying a face
of a cube with a hexad or a vertex with an octad.

\subsection{Strata Correspondence and Positivity of the Triple Freedom}

The five-region structure of channel \(\pi\) has types
\[
 1_\perp+1_{--}+3_{++}.
\]
The incidence on \(B_K\) has types
\[
 1_{\rm transversal}+1_{H^-}+3_{\rm restantes}
\]
after marking the negative hexad. There is therefore a natural
correspondence between
\emph{role} strata:
\begin{equation}
 1_\perp\longleftrightarrow O_{B_K}^\perp,
 \qquad
 1_{--}\longleftrightarrow
 \widehat\iota_{D,\ell}(H_\alpha^-),
 \qquad
 3_{++}\longleftrightarrow
 \mathcal O_{B_K}^{+,3}.
 \label{p12-83-c20-s3:eq:correspondencia-estratos-cinco-cinco}
\end{equation}
Here \(\mathcal O_{B_K}^{+,3}\) is the unordered set of the other
three residual octads. An individual bijection between the three positive
regions and those three octads requires declaring an additional gauge of
\(S_3\). The coincidence of the pattern \(1+1+3\) does not by itself determine that
ordering.

\subsection{Automorphism group of the octad system}

\begin{theorem}[Mathieu group \(M_{24}\)]
\label{p12-83-c20-s3:thm:automorfismos-w24-m24}
The automorphism group of the octad Steiner system is
\begin{equation}
 \boxed{
 \operatorname{Aut}(W_{24})\cong M_{24},
 \qquad
 |M_{24}|=244823040.}
 \label{p12-83-c20-s3:eq:automorfismos-w24-m24}
\end{equation}
\end{theorem}

The theorem identifies the classical automorphism group of the input
binary data. It does not assert that \(M_{24}\) is generated from a single
dodecad, or that the ternary code \(C_W\) determines
\(\mathcal G_{24}\) by itself.

\subsection{Categorical separation from the lattice branch}

The branch constructed in this unit is
\begin{equation}
 (W_{12}^{\rm HMT};\mathcal G_{24},D,\ell)
 \longrightarrow W_{24}
 \longrightarrow\operatorname{Aut}(W_{24})\cong M_{24}.
 \label{p12-83-c20-s3:eq:rama-binaria-w24-m24}
\end{equation}
The lattice branch to be constructed subsequently is
\begin{equation}
 C_W\longrightarrow N(A_2^{12})
 \xrightarrow{\text{index-three neighbor}}\Lambda_{24}.
 \label{p12-83-c20-s3:eq:rama-ternaria-niemeier-leech-anunciada}
\end{equation}
Their constructors, domains, and codomains are distinct. In particular,
\begin{equation}
 \boxed{
 \text{no arrow is introduced }
 W_{24}\longrightarrow\Lambda_{24}.}
 \label{p12-83-c20-s3:eq:no-flecha-w24-leech}
\end{equation}

\begin{figure}[htbp]
\centering
\begin{tikzpicture}[
  >=Stealth,
  node distance=13mm and 16mm,
  box/.style={draw,rounded corners=1.5mm,align=center,
    inner xsep=3mm,inner ysep=2.2mm},
  arr/.style={->,thick}
]
\node[box] (cw) {\(C_W\)};
\node[box,right=of cw] (w12) {\(W_{12}\)};
\node[box,right=of w12] (w24) {\(W_{24}\)};
\node[box,right=of w24] (m24) {\(M_{24}\)};
\node[box,below=of cw] (cw2) {\(C_W\)};
\node[box,right=of cw2] (n) {\(N(A_2^{12})\)};
\node[box,right=of n] (leech) {\(\Lambda_{24}\)};
\draw[arr] (cw)--(w12);
\draw[arr] (w12)--node[above,font=\scriptsize]
 {\((\mathcal G_{24},D,\ell)\)} (w24);
\draw[arr] (w24)--(m24);
\draw[arr] (cw2)--node[below,font=\scriptsize]
 {pegado ternario} (n);
\draw[arr] (n)--node[below,font=\scriptsize]
 {\(3\)-neighbor} (leech);
\node[below=5mm of w24,font=\footnotesize,align=center]
 {Vertical alignment is not an arrow.};
\end{tikzpicture}
\caption{Typed bifurcation of the binary and ternary extensions.}
\label{p12-83-c20-s3:fig:bifurcacion-w24-leech-autonoma}
\end{figure}

\subsection{Obstructions to the binary prolongation \(W_{12}\to W_{24}\): Golay code, octads, alignment and separation of the Leech branch}

\begin{criterion}[Refutation of the binary prolongation]
The construction is refuted if any of the following facts occurs:
\begin{enumerate}
 \item the weight spectrum of \(\mathcal G_{24}\) contains a weight not
       listed in \eqref{p12-83-c20-s3:eq:pesos-golay-binario-extendido};
 \item the supports of weight eight not form \(S(5,8,24)\);
 \item a residual intersection of an octad with \(D\) that contains five
       points of \(D\) has cardinality other than six;
 \item a residual hexad admits either no lifting octad or more than one;
 \item \(\iota_D\) is applied directly to an HMT hexad without the alignment \(\ell\);
 \item the incidence law on a tetrad is not \(4+1\);
 \item the positive triple is ordered without declaring the gauge \(S_3\);
 \item \(\operatorname{Aut}(W_{24})\) has no order \(244823040\);
 \item it is asserted that \(W_{24}\) arises from \(W_{12}\) without the binary datum
       \(\mathcal G_{24}\), the dodecad, and the alignment;
 \item the two branches are serialized by means of an arrow
       \(W_{24}\to\Lambda_{24}\).
\end{enumerate}
\end{criterion}


% Unidad U021. Pegado discriminante ternario y Niemeier A2^12.
% Redaccion autonoma desde la autoridad material 1063.

\section{Ternary discriminant gluing and construction of the Niemeier lattice of type \texorpdfstring{\(A_2^{12}\)}{A₂¹²}}
\label{p12-83-c20-s4:chap:pegado-niemeier-a2-doce}

This unit begins the lattice branch directly from the self-dual ternary code
\(C_W\). The binary system \(W_{24}\), the dodecad, and the group
\(M_{24}\) are not inputs of this construction. The typed sequence is
\[
 C_W
 \longrightarrow
 (A_2^*/A_2)^{12}
 \longrightarrow
 N(C_W)
 \cong N(A_2^{12}).
\]

\subsection{Reticulo root \texorpdfstring{\(A_2\)}{A₂}}

Let
\[
 A_2=\mathbb Z\alpha_1\oplus\mathbb Z\alpha_2
\]
the root lattice whose Gram matrix, in the ordered basis
\((\alpha_1,\alpha_2)\), is
\begin{equation}
 G_{A_2}=
 \begin{pmatrix}
 2&-1\\
 -1&2
 \end{pmatrix}.
 \label{p12-83-c20-s4:eq:gram-a2-autonomo}
\end{equation}
For \(x=x_1\alpha_1+x_2\alpha_2\),
\begin{equation}
 \langle x,x\rangle
 =2(x_1^2-x_1x_2+x_2^2).
 \label{p12-83-c20-s4:eq:norma-a2-autonoma}
\end{equation}
In particular, \(A_2\) is even, positive defined and
\[
 \det A_2=3.
\]

We define
\[
 \rho:=\alpha_1+\alpha_2.
\]
Then
\begin{equation}
 \rho^2=2,
 \qquad
 \langle\rho,\alpha_1\rangle
 =\langle\rho,\alpha_2\rangle=1.
 \label{p12-83-c20-s4:eq:identidades-rho-a2}
\end{equation}

\subsection{Discriminant module and form}

The dual lattice is
\[
 A_2^*
 :=\{x\in A_2\otimes_{\mathbb Z}\mathbb Q:
       \langle x,A_2\rangle\subseteq\mathbb Z\}.
\]
We choose the representatives
\begin{equation}
 \mu_0=0,
 \qquad
 \mu_1=\frac{2\alpha_1+\alpha_2}{3},
 \qquad
 \mu_2=\frac{\alpha_1+2\alpha_2}{3}.
 \label{p12-83-c20-s4:eq:representantes-discriminantes-a2}
\end{equation}
Their pairings with the simple base are
\[
\begin{array}{c|cc}
 &\alpha_1&\alpha_2\\ \hline
 \mu_1&1&0\\
 \mu_2&0&1
\end{array},
\]
thereby \(\mu_1,\mu_2\in A_2^*\). Furthermore,
\[
 3\mu_1=2\alpha_1+\alpha_2\in A_2,
 \qquad
 3\mu_2=\alpha_1+2\alpha_2\in A_2,
\]
while \(\mu_1,\mu_2\notin A_2\). It is obtained
\begin{equation}
 A_2^*/A_2
 =\{\overline\mu_0,\overline\mu_1,\overline\mu_2\}
 \cong\mathbb Z/3\mathbb Z
 \cong\mathbb F_3.
 \label{p12-83-c20-s4:eq:discriminante-a2-f3}
\end{equation}

The sum of representatives is explicitly reduced by means of
\[
 \mu_1+\mu_1-\mu_2=\alpha_1,
 \qquad
 \mu_1+\mu_2=\rho,
 \qquad
 \mu_2+\mu_2-\mu_1=\alpha_2.
\]
Consequently,
\begin{equation}
 \mu_r+\mu_s-\mu_{r+s\,({\rm mod}\,3)}\in A_2
 \qquad(r,s\in\mathbb F_3).
 \label{p12-83-c20-s4:eq:adicion-representantes-discriminantes-a2}
\end{equation}

\begin{proposition}[Discriminant Form of \(A_2\)]
\label{p12-83-c20-s4:prop:forma-discriminante-a2}
Under \eqref{p12-83-c20-s4:eq:discriminante-a2-f3}, the discriminant bilinear and quadratic
forms are
\begin{align}
 b(\overline\mu_r,\overline\mu_s)
 &\equiv\frac{2rs}{3}\pmod{\mathbb Z},
 \label{p12-83-c20-s4:eq:bilineal-discriminante-a2}\\
 q(\overline\mu_r)
 &\equiv\frac{2r^2}{3}\pmod{2\mathbb Z}.
 \label{p12-83-c20-s4:eq:cuadratica-discriminante-a2}
\end{align}
\end{proposition}

\begin{proof}
The formula \eqref{p12-83-c20-s4:eq:norma-a2-autonoma} gives
\[
 \mu_1^2=\mu_2^2=\frac23,
 \qquad
 \langle\mu_1,\mu_2\rangle=\frac13.
\]
For \(r,s\in\{0,1,2\}\), these values coincide, modulo
\(\mathbb Z\), with \(2rs/3\). The quadratic formula results by taking
\(r=s\) and recalling that the lattice is even.
\end{proof}

\subsection{\(12\) factors and gluing map}

Let
\[
 L_0:=A_2^{12}.
\]
Then
\[
 L_0^*=(A_2^*)^{12},
 \qquad
 L_0^*/L_0\cong\mathbb F_3^{12}.
\]
For \(c=(c_1,\ldots,c_{12})\in\mathbb F_3^{12}\), define
\begin{equation}
 \phi(c):=(\mu_{c_1},\ldots,\mu_{c_{12}})\in L_0^*.
 \label{p12-83-c20-s4:eq:aplicacion-pegado-cw}
\end{equation}
The extension determined by \(C_W\) is
\begin{equation}
 \begin{aligned}
 N(C_W)
 &:=\bigcup_{c\in C_W}(L_0+\phi(c))\\
 &=\{x\in L_0^*:x\bmod L_0\in C_W\}.
 \end{aligned}
 \label{p12-83-c20-s4:eq:definicion-niemeier-cw}
\end{equation}

\begin{lemma}[Additive closure of the gluing]
\label{p12-83-c20-s4:lem:clausura-pegado-cw}
The set \(N(C_W)\) is a discrete subgroup of \(L_0^*\) of rank
twenty-four.
\end{lemma}

\begin{proof}
Let \(x=\ell+\phi(c)\) and \(y=m+\phi(d)\), with \(\ell,m\in L_0\) and
\(c,d\in C_W\). By \eqref{p12-83-c20-s4:eq:adicion-representantes-discriminantes-a2},
\[
 \phi(c)+\phi(d)-\phi(c+d)\in L_0.
\]
Since \(C_W\) is linear, \(c+d\in C_W\), and then
\(x+y\in L_0+\phi(c+d)\). The same identity applied to the opposite gives
inverses. The inclusions
\[
 L_0\subseteq N(C_W)\subseteq L_0^*
\]
They test discretion and range twenty-four.
\end{proof}

\subsection{Integrality induced by self-duality}

The discriminant bilinear form of \(L_0\) is the orthogonal sum of twelve
copies of \eqref{p12-83-c20-s4:eq:bilineal-discriminante-a2}. For
\(c,d\in\mathbb F_3^{12}\),
\begin{equation}
 b(c,d)
 \equiv\frac23\sum_{i=1}^{12}c_id_i
 \pmod{\mathbb Z}.
 \label{p12-83-c20-s4:eq:bilineal-discriminante-a2-doce}
\end{equation}

\begin{lemma}[Integrity of the gluing]
\label{p12-83-c20-s4:lem:integralidad-pegado-cw}
The lattice \(N(C_W)\) is integral.
\end{lemma}

\begin{proof}
The self-duality \(C_W=C_W^\perp\) implica
\[
 \sum_{i=1}^{12}c_id_i\equiv0\pmod3
 \qquad(c,d\in C_W).
\]
By \eqref{p12-83-c20-s4:eq:bilineal-discriminante-a2-doce}, the discriminant form
vanishes on \(C_W\times C_W\). Therefore, the pairing of any
two glued classes is integral.
\end{proof}

\subsection{Parity Induced by the Code Weights}

At a nonzero coordinate of the code, the discriminant quadratic form
equals \(2/3\) modulo \(2\mathbb Z\). Therefore,
\begin{equation}
 q(c)\equiv\frac23\operatorname{wt}(c)\pmod{2\mathbb Z}.
 \label{p12-83-c20-s4:eq:cuadratica-discriminante-cw}
\end{equation}

\begin{lemma}[Parity of the gluing]
\label{p12-83-c20-s4:lem:paridad-pegado-cw}
The lattice \(N(C_W)\) is even.
\end{lemma}

\begin{proof}
The only weights of \(C_W\) are \(0,6,9,12\), and
\[
 \frac23\{0,6,9,12\}=\{0,4,6,8\}\subseteq2\mathbb Z.
\]
Every glue class is isotropic for the discriminant quadratic form, so all norms of \(N(C_W)\) are even.
\end{proof}

Let \(g_1,\ldots,g_6\) be the rows of \(G_W=(I_6\mid A_W)\) and
\(\gamma_i=\phi(g_i)\). Exact multiplication in \(A_2^*\) gives
\begin{equation}
 (\langle\gamma_i,\gamma_j\rangle)_{i,j=1}^{6}
 =
 \begin{pmatrix}
 4&2&2&2&2&2\\
 2&4&2&2&2&2\\
 2&2&4&2&2&2\\
 2&2&2&4&2&2\\
 2&2&2&2&4&2\\
 2&2&2&2&2&4
 \end{pmatrix}.
 \label{p12-83-c20-s4:eq:gram-generadores-pegado-cw}
\end{equation}
This matrix provides an additional direct control of integrality and parity for the six gluing generators.

\subsection{Index, determinant, and unimodularity}

The cosets of \(L_0\) that appear in
\eqref{p12-83-c20-s4:eq:definicion-niemeier-cw} are indexed by the
\(|C_W|=729=3^6\) elements of the code and are distinct. Consequently,
\begin{equation}
 [N(C_W):L_0]=3^6.
 \label{p12-83-c20-s4:eq:indice-niemeier-sobre-a2-doce}
\end{equation}
Since
\[
 \det L_0=(\det A_2)^{12}=3^{12},
\]
the determinant formula for a finite extension gives
\begin{equation}
 \det N(C_W)
 =\frac{\det L_0}{[N(C_W):L_0]^2}
 =\frac{3^{12}}{(3^6)^2}=1.
 \label{p12-83-c20-s4:eq:determinante-niemeier-cw}
\end{equation}

\subsection{Root system induced by Witt incidence and classification of its orbits}

En each one of the two classes discriminantes not nulas of
\(A_2^*/A_2\), the norm minima is
\[
 \mu_1^2=\mu_2^2=\frac23.
\]
If \(c\in C_W\setminus\{0\}\), then
\(\operatorname{wt}(c)\geq6\). By orthogonality of the twelve factors,
any vector of \(L_0+\phi(c)\) has norm at least
\begin{equation}
 6\cdot\frac23=4.
 \label{p12-83-c20-s4:eq:cota-clases-no-triviales-niemeier}
\end{equation}
Therefore, no nontrivial gluing class contains vectors of norm
two. The roots of \(N(C_W)\) are exactly the roots already present in
\(L_0=A_2^{12}\).

\begin{theorem}[Realization of the Niemeier of type \(A_2^{12}\)]
\label{p12-83-c20-s4:thm:niemeier-a2-doce-autonomo}
The lattice \(N(C_W)\) is positive definite, even, and unimodular, and its root
system is \(A_2^{12}\). Consequently,
\begin{equation}
 \boxed{N(C_W)\cong N(A_2^{12}).}
 \label{p12-83-c20-s4:eq:identificacion-niemeier-a2-doce}
\end{equation}
\end{theorem}

\begin{proof}
Positivity follows from \(L_0^*\). The
\cref{p12-83-c20-s4:lem:integralidad-pegado-cw,p12-83-c20-s4:lem:paridad-pegado-cw} proves
integrality and parity; \eqref{p12-83-c20-s4:eq:determinante-niemeier-cw} proves
unimodularity. The bound \eqref{p12-83-c20-s4:eq:cota-clases-no-triviales-niemeier}
identifies the root system. The Niemeier classification uniquely associates
the indicated lattice with that root system.
\end{proof}

\subsection{Obstructions to the ternary gluing \(C_W\to N(A_2^{12})\): parity, unimodularity, and root system}

\begin{criterion}[Refutation of ternary gluing]
The construction is refuted if any of the following facts occurs:
\begin{enumerate}
 \item any of the representatives \(\mu_0,\mu_1,\mu_2\) does not belong to
       \(A_2^*\), two of its classes coincide, or
       \eqref{p12-83-c20-s4:eq:adicion-representantes-discriminantes-a2} fails;
 \item the discriminant form is not the one given in
       \cref{p12-83-c20-s4:prop:forma-discriminante-a2};
 \item \(N(C_W)\) is not closed under sum or inversos;
 \item there exists \(c,d\in C_W\) with
       \(\frac23\sum_i c_id_i\notin\mathbb Z\);
 \item one palabra of \(C_W\) produces
       \(\frac23\operatorname{wt}(c)\notin2\mathbb Z\);
 \item the index is not \(3^6\) or the determinant is not one;
 \item a nontrivial gluing class contains a vector of norm two;
 \item \(W_{24}\), a dodecad, or a binary datum is used in any of
the preceding gluing operations.
\end{enumerate}
\end{criterion}


% Unidad U022. Origen incidencial, vecino ternario y Leech.
% Redaccion autonoma desde la autoridad material 1063.

\section{Incidental origin, typed radial problem and construction of the rootless unimodular neighbor}
\label{p12-83-c20-s5:chap:origen-incidencial-vecino-leech}

The lattice \(N(A_2^{12})\) has already been constructed. This unit determines a
base point from an ordered incidence frame, solves a
variational problem on a declared domain, verifies the three conditions
for the index-three neighbor, and proves that the result is the Leech
lattice. The neighbor vector is denoted \(v_{\mathcal F_*}\): it depends on the
incidence frame and does not enter the numerical carry of \(\alpha\).

\subsection{Ordered set of incidences}

Consider the four hexads
\begin{align}
 H^-&=\{4,5,6,7,9,10\},
 \label{p12-83-c20-s5:eq:hexada-marco-negativa}\\
 H_e&=\{1,3,4,6,9,10\},
 \label{p12-83-c20-s5:eq:hexada-marco-e}\\
 H_\varphi&=\{1,2,3,4,7,9\},
 \label{p12-83-c20-s5:eq:hexada-marco-phi}\\
 H_{\rm APP}&=\{2,3,5,10,11,12\}.
 \label{p12-83-c20-s5:eq:hexada-marco-app}
\end{align}
The datum is the ordered framework.
\begin{equation}
 \mathcal F:=(H^-,H_e,H_\varphi,H_{\rm APP}).
 \label{p12-83-c20-s5:eq:marco-incidencial-ordenado}
\end{equation}
We define
\begin{equation}
 o(\mathcal F)
 :=(H_e\cap H_\varphi)\setminus(H^-\cup H_{\rm APP}).
 \label{p12-83-c20-s5:eq:operador-origen-incidencial}
\end{equation}

\begin{theorem}[Equivariant Reconstruction of the Origin]
\label{p12-83-c20-s5:thm:origen-incidencial-equivariante}
For the frame canonical \(\mathcal F_*\),
\begin{equation}
 \boxed{o(\mathcal F_*)=\{1\}.}
 \label{p12-83-c20-s5:eq:origen-incidencial-singleton}
\end{equation}
Moreover, for every simultaneous permutation \(g\in M_{12}\),
\begin{equation}
 o(g\mathcal F)=g\,o(\mathcal F).
 \label{p12-83-c20-s5:eq:equivarianza-origen-incidencial}
\end{equation}
\end{theorem}

\begin{proof}
We have
\[
 H_e\cap H_\varphi=\{1,3,4,9\}
\]
and
\[
 (H^-\cup H_{\rm APP})\cap\{1,3,4,9\}=\{3,4,9\}.
\]
The difference is \(\{1\}\). For equivariance one applies the
identities
\[
 g(A\cap B)=gA\cap gB,
 \quad
 g(A\cup B)=gA\cup gB,
 \quad
 g(A\setminus B)=gA\setminus gB,
\]
valid for all bijections.
\end{proof}

The tetrad \(B_K=\{4,6,9,10\}\) does not by itself determine the origin. The
singleton is obtained from the complete ordered frame; it is not a
coordinate chosen after the result has been observed.

\subsection{Positive radial chamber of the root system and selection of the dominant cone}

Let \(\rho_i\) the copia of \(\rho=\alpha_1+\alpha_2\) in the
\(i\)-esimo factor of \(A_2^{12}\). Consideramos the dominio
\begin{equation}
 \mathscr R_+
 :=\left\{
 v(a)=\sum_{i=1}^{12}a_i\rho_i:
 a_i=1+3b_i, b_i\in\mathbb Z_{\geq0}
 \right\}.
 \label{p12-83-c20-s5:eq:camara-radial-positiva}
\end{equation}
Since the factors are orthogonal and \(\rho_i^2=2\), if
\(S=\sum_i b_i\), then
\begin{align}
 v(a)^2
 &=2\sum_{i=1}^{12}(1+3b_i)^2\notag\\
 &=24+12S+18\sum_{i=1}^{12}b_i^2.
 \label{p12-83-c20-s5:eq:norma-radial-expandida}
\end{align}
Reducing modulo eighteen,
\[
 v(a)^2\equiv6+12S\pmod{18},
\]
and therefore
\begin{equation}
 v(a)^2\equiv0\pmod{18}
 \quad\Longleftrightarrow\quad
 S\equiv1\pmod3.
 \label{p12-83-c20-s5:eq:congruencia-radial-suma-b}
\end{equation}

\begin{proposition}[Minima in the declared radial chamber]
\label{p12-83-c20-s5:prop:minimos-camara-radial}
On the subset of \(\mathscr R_+\) defined by
\(v(a)^2\equiv0\pmod{18}\), the minimum norm is \(54\). It is attained
exactly at twelve vectors: for each \(j\in\{1,\ldots,12\}\), the
coefficient \(a_j\) equals four, and the remaining eleven equal one.
\end{proposition}

\begin{proof}
Condition \eqref{p12-83-c20-s5:eq:congruencia-radial-suma-b} and the nonnegativity of
the \(b_i\) imply that the smallest possible value of \(S\) is one. This
forces a unique index \(j\) with \(b_j=1\) and \(b_i=0\) for
\(i\neq j\). Substituting into \eqref{p12-83-c20-s5:eq:norma-radial-expandida},
\[
 v(a)^2=24+12+18=54.
\]
There are twelve choices of \(j\). The next admissible value of \(S\) is
four and produces a strictly greater norm.
\end{proof}

The origin \(o(\mathcal F_*)=\{1\}\) selects
\begin{equation}
 \boxed{
 v_{\mathcal F_*}
 =4\rho_1+\rho_2+\cdots+\rho_{12},
 \qquad
 v_{\mathcal F_*}^2=54.}
 \label{p12-83-c20-s5:eq:vector-marco-incidencial}
\end{equation}

\begin{remark}[Exact scope of minimality]
The norm \(54\) is minimal in the positive radial chamber
\eqref{p12-83-c20-s5:eq:camara-radial-positiva}, not in the entire residue class. The vector
\[
 u=(\alpha_1-2\alpha_2,\rho,\ldots,\rho)
\]
satisfies
\[
 \alpha_1-2\alpha_2-\rho=-3\alpha_2,
\]
so that
\[
 u\equiv(\rho,\ldots,\rho)
 \equiv v_{\mathcal F_*}\pmod{3A_2^{12}}.
\]
However,
\[
 (\alpha_1-2\alpha_2)^2
 =2(1+2+4)=14,
 \qquad
 u^2=14+11\cdot2=36.
\]
This witness prevents extending the minimality assertion outside the domain
of \cref{p12-83-c20-s5:prop:minimos-camara-radial}.
\end{remark}

\subsection{Character and congruence kernel}

We write
\[
 N:=N(A_2^{12}),
 \qquad
 v:=v_{\mathcal F_*}.
\]
We define
\begin{equation}
 \chi_v:N\longrightarrow\mathbb Z/3\mathbb Z,
 \qquad
 \chi_v(x)=\langle x,v\rangle\bmod3,
 \label{p12-83-c20-s5:eq:caracter-vecino-leech}
\end{equation}
its nucleus
\begin{equation}
 N_0:=\ker\chi_v
 =\{x\in N:\langle x,v\rangle\equiv0\pmod3\},
 \label{p12-83-c20-s5:eq:nucleo-congruencia-vecino}
\end{equation}
and the extension
\begin{equation}
 N_v:=N_0+\mathbb Z\frac v3.
 \label{p12-83-c20-s5:eq:definicion-vecino-indice-tres}
\end{equation}

\subsection{Condition I1: surjectivity and exclusion of roots}

By construction, \(v\in A_2^{12}\subseteq N\). The roots of a factor
\(A_2\) are
\[
 \pm\alpha_1,
 \quad
 \pm\alpha_2,
 \quad
 \pm\rho.
\]
Their pairings with \(\rho\) are \(\pm1\) or \(\pm2\). All the
coefficients of \(v\) are congruent to one modulo three:
\[
 4\equiv1\pmod3,
 \qquad
 1\equiv1\pmod3.
\]
Therefore, no root of \(A_2^{12}\) has zero pairing with
\(v\) modulo three. In particular,
\[
 \chi_v(\alpha_{1,1})=1,
\]
so that \(\chi_v\) is surjective and
\begin{equation}
 [N:N_0]=3.
 \label{p12-83-c20-s5:eq:indice-nucleo-vecino}
\end{equation}
Moreover, no root of \(N\) belongs to \(N_0\).

\subsection{Condition I2: norm divisibility}

The identity \eqref{p12-83-c20-s5:eq:vector-marco-incidencial} gives
\begin{equation}
 v^2=54\equiv0\pmod{18},
 \qquad
 \left(\frac v3\right)^2=6.
 \label{p12-83-c20-s5:eq:condicion-i2-vecino}
\end{equation}
In particular, \(\langle v,v\rangle\equiv0\pmod3\), hence
\(v\in N_0\).

\subsection{Condition I3: class dual of order \(3\)}

If \(x\in N_0\), there exists \(k\in\mathbb Z\) with
\(\langle x,v\rangle=3k\). Therefore,
\[
 \left\langle x,\frac v3\right\rangle=k\in\mathbb Z,
\]
and
\begin{equation}
 \frac v3\in N_0^*.
 \label{p12-83-c20-s5:eq:v-tercio-en-dual-nucleo}
\end{equation}
On the other hand,
\[
 \left\langle\frac v3,\alpha_{1,1}\right\rangle=\frac43.
\]
Since \(\alpha_{1,1}\in N\) and \(N\) is integral, this proves
\begin{equation}
 \frac v3\notin N.
 \label{p12-83-c20-s5:eq:v-tercio-no-en-niemeier}
\end{equation}
Finally, \(3(v/3)=v\in N_0\). The class of \(v/3\) in
\(N_0^*/N_0\) is nontrivial and has order exactly three. Consequently,
\begin{equation}
 [N_v:N_0]=3.
 \label{p12-83-c20-s5:eq:indice-extension-vecino}
\end{equation}

\subsection{Parity, integrality and determinant of the neighbor}

\begin{proposition}[Unimodular structure of the neighbor]
\label{p12-83-c20-s5:prop:estructura-unimodular-vecino}
The lattice \(N_v\) is integral, even, and unimodular.
\end{proposition}

\begin{proof}
Since \(N\) is unimodular and \([N:N_0]=3\),
\[
 \det N_0=[N:N_0]^2\det N=9.
\]
Since \([N_v:N_0]=3\),
\[
 \det N_v=\frac{\det N_0}{[N_v:N_0]^2}=1.
\]

Let \(y=x+m(v/3)\), with \(x\in N_0\), and write
\(\langle x,v\rangle=3k\). Then
\begin{align*}
 y^2
 &=x^2+2m\left\langle x,\frac v3\right\rangle
       +m^2\left(\frac v3\right)^2\\
 &=x^2+2mk+6m^2.
\end{align*}
The three summands form an even integer. The polarized form proves
integrality, and determinant one proves unimodularity.
\end{proof}

\subsection{\(2\) clasis globalis and \(6\) clasis local}

Globally, \eqref{p12-83-c20-s5:eq:definicion-vecino-indice-tres} adds only two nontrivial
cosets:
\[
 N_0+\frac v3,
 \qquad
 N_0-\frac v3.
\]
Locally, in each factor \(A_2\), a coordinate of those classes belongs to one of the six classes
\begin{equation}
 \mathscr C_{j,\pm}
 :=A_2+\mu_j\pm\frac\rho3,
 \qquad j\in\{0,1,2\}.
 \label{p12-83-c20-s5:eq:seis-clases-locales-vecino}
\end{equation}
In the first coordinate, \(4\rho/3\equiv\rho/3\pmod{A_2}\), so
the same list holds for the twelve factors.

\begin{table}[htbp]
\centering
\caption{Minimal representatives of the six local cosets.}
\label{p12-83-c20-s5:tab:seis-clases-locales-vecino}
\begin{tabular}{cclc}
\toprule
class&minimal representative&\((u,v)\)&
\(2(u^2-uv+v^2)/9\)\\
\midrule
\(\mathscr C_{0,-}\)&\(-\rho/3\)&\((-1,-1)\)&\(2/9\)\\
\(\mathscr C_{0,+}\)&\( \rho/3\)&\((1,1)\)&\(2/9\)\\
\(\mathscr C_{1,-}\)&\( \mu_1-\rho/3\)&\((1,0)\)&\(2/9\)\\
\(\mathscr C_{1,+}\)&\( \mu_1+\rho/3-\rho\)&\((0,-1)\)&\(2/9\)\\
\(\mathscr C_{2,-}\)&\( \mu_2-\rho/3\)&\((0,1)\)&\(2/9\)\\
\(\mathscr C_{2,+}\)&\( \mu_2+\rho/3-\rho\)&\((-1,0)\)&\(2/9\)\\
\bottomrule
\end{tabular}
\end{table}

Here \((u,v)\) means the representative
\((u\alpha_1+v\alpha_2)/3\).

\begin{lemma}[Exact minimum of the six classes]
\label{p12-83-c20-s5:lem:minimo-seis-clases-locales}
Each class in \eqref{p12-83-c20-s5:eq:seis-clases-locales-vecino} has minimum norm
exactly \(2/9\).
\end{lemma}

\begin{proof}
The norm of \((u\alpha_1+v\alpha_2)/3\) is
\[
 \frac29(u^2-uv+v^2).
\]
The integral form \(u^2-uv+v^2\) is positive definite. On a nonzero
residual class it does not take the value zero, and hence is at least one. The six
representatives in \cref{p12-83-c20-s5:tab:seis-clases-locales-vecino} attain one.
\end{proof}

\subsection{Criterion for the absence of roots in the Leech lattice and consequence for the minimum distance}

\begin{theorem}[Absence of vectors of norm two]
\label{p12-83-c20-s5:thm:vecino-sin-raices}
The lattice \(N_v\) contains no vectors of norm two.
\end{theorem}

\begin{proof}
All roots of \(N\) are those of \(A_2^{12}\). Condition I1 proves
that none belongs to \(N_0\). Hence class \(N_0\) contains no
vectors of norm two.

In a new class \(N_0\pm v/3\), each of the twelve coordinates
belongs to one of the six local classes. By
\cref{p12-83-c20-s5:lem:minimo-seis-clases-locales}, each coordinate contributes at least
\(2/9\). Orthogonality gives
\[
 \left\|x\pm\frac v3\right\|^2
 \geq12\cdot\frac29=\frac83>2.
\]
Nor do the new classes contain roots.
\end{proof}

\subsection{Identification with the Leech lattice}

\begin{theorem}[Construction of the Leech lattice]
\label{p12-83-c20-s5:thm:identificacion-vecino-leech}
The lattice
\begin{equation}
 N_v
 =\ker(\langle\,\cdot\,,v_{\mathcal F_*}\rangle\bmod3)
  +\mathbb Z\frac{v_{\mathcal F_*}}3
 \label{p12-83-c20-s5:eq:formula-final-vecino-leech}
\end{equation}
is isometric to the Leech lattice:
\begin{equation}
 \boxed{N_v\cong\Lambda_{24}.}
 \label{p12-83-c20-s5:eq:identificacion-leech-autonoma}
\end{equation}
\end{theorem}

\begin{proof}
The lattice \(N_v\) has rank twenty-four and is positive definite.
\cref{p12-83-c20-s5:prop:estructura-unimodular-vecino} proves that it is even and unimodular;
\cref{p12-83-c20-s5:thm:vecino-sin-raices} proves that it has no roots. By the
Niemeier classification, there exists a unique isometry class of positive-definite,
even, unimodular lattices of rank twenty-four without roots:
the Leech lattice.
\end{proof}

\subsection{The 243 syndromes of the ternary punctured code}

Let \(G_{11}\) be the ternary Golay code with parameters
\([11,6,5]_3\). Its syndrome space has cardinality
\begin{equation}
 3^{11-6}=3^5=243.
 \label{p12-83-c20-s5:eq:numero-sindromes-golay-ternario}
\end{equation}
The number of weight errors of at most two is
\begin{equation}
 1+2\binom{11}{1}+2^2\binom{11}{2}
 =1+22+220=243.
 \label{p12-83-c20-s5:eq:bola-radio-dos-golay-ternario}
\end{equation}

\begin{proposition}[Perfect correspondence error--syndrome]
\label{p12-83-c20-s5:prop:correspondencia-error-sindrome-golay}
Each syndrome of \(G_{11}\) has a unique error representative of weight at most two.
\end{proposition}

\begin{proof}
Minimum distance five implies that two distinct errors of weight at
most two cannot differ by a codeword: their difference would have
weight at most four. The error--syndrome map is therefore
injective on the ball of radius two. Both sets have cardinality
\(243\) by \eqref{p12-83-c20-s5:eq:numero-sindromes-golay-ternario} and
\eqref{p12-83-c20-s5:eq:bola-radio-dos-golay-ternario}; the map is bijective.
\end{proof}

These \(243\) syndromes are not the \(243\) visible words obtained from the
TPK catalogue. They share a cardinal, but they possess distinct domains, equivalences and
functions.

\subsection{Obstructions of the ternary neighbor \(N_v\cong\Lambda_{24}\): class of order \(3\), absence of roots, and separation of the cardinals \(243\)}

\begin{criterion}[Refutation of the neighbor and of the Leech identification]
The construction is refuted if any of the following facts occurs:
\begin{enumerate}
 \item the operator \(o(\mathcal F)\) not produces \(\{1\}\) or is not
       equivariant;
 \item there exists in the radial chamber an admissible vector of norm less than
       \(54\), or the number of minimizers is not twelve;
 \item global minimality is asserted despite the witness of norm \(36\);
 \item \(\chi_v\) is not sobreyectivo or one root belongs to \(N_0\);
 \item \(v^2\not\equiv0\pmod{18}\) o \((v/3)^2\neq6\);
 \item \(v/3\notin N_0^*\), \(v/3\in N\), or its class does not have order three;
 \item the determinant of \(N_v\) is not one, or the neighbor ceases to be even;
 \item some of the six clasis local has minimum least that \(2/9\);
 \item one of the two new global classes contains a root;
 \item the two global classes are confused with the six local classes;
 \item \(W_{24}\) is used as an antecedent of the gluing or of the neighbor;
 \item the \(243\) syndromes are identified with the \(243\) words
       visible by sharing cardinality.
\end{enumerate}
\end{criterion}


\section{Finite holonomy, Schläfli configuration and symmetry \texorpdfstring{\(E_6\)}{E6}}
\label{p12-83-c20-s6:chap:holonomia-schlafli-e6-nuevo}
\index{finite holonomy}
\index{Schläfli configuration}
\index{E6@\(E_6\)}

The Witt branch does not exhaust the exceptional structures accessible from
APP. The parallel branch of this chapter starts from the two mixed
sum--product and product--sum polarizations, together with a plaquette
orientation and a declared central section. The calculated curvature determines
the alternating form; this defines the central extension. An
extension is not inferred from cardinality 27 alone.

\subsection{Mixed APP curvatures}

On \(T_9=(\mathbb Z/9\mathbb Z)^2\), with
\(S(x,y)=x+y\) and \(P(x,y)=xy\), let
\[
 \Delta_xf=f(x+1,y)-f(x,y),\qquad
 \Delta_yf=f(x,y+1)-f(x,y).
\]
Mixed connections are
\[
 A^{SP}=(\Delta_xS,\Delta_yP)=(1,x),\qquad
 A^{PS}=(\Delta_xP,\Delta_yS)=(y,1).
\]
Its discrete curvature is
\[
 F_A=\Delta_xA_y-\Delta_yA_x.
\]

\begin{proposition}[Oriented opposite curvatures]
\label{p12-83-c20-s6:prop:curvaturas-app-e6-nuevo}
\[
 F_{A^{SP}}=1,\qquad F_{A^{PS}}=-1\pmod9.
\]
The circulation over an oriented rectangle of sides \(r,s\) is,
respectively, \(rs\) and \(-rs\).
\end{proposition}

\begin{proof}
\(\Delta_xx-\Delta_y1=1\) and
\(\Delta_x1-\Delta_yy=-1\). A rectangle contains \(rs\) plaquettes; upon
summing, the interior edges cancel and the boundary circulation remains.
\end{proof}

For displazamientos \(u=(a,b)\), \(v=(c,d)\), the class alternante is
\begin{equation}
 \sigma(u,v)=ad-bc\pmod9.
 \label{p12-83-c20-s6:eq:forma-alternante-e6-nueva}
\end{equation}
The orientation fixes the sign of \(\sigma\). The central section used
below belongs to the gauge of this realization; replacing it with a coboundary
produces an isomorphic extension, not a new invariant.

\subsection{Central extension and ternary reduction}

We define the nonadic Heisenberg group
\begin{equation}
 \mathcal H_9=\{(a,b,z):a,b,z\in\mathbb Z/9\mathbb Z\},
 \quad
 (a,b,z)(c,d,w)=(a+c,b+d,z+w+bc).
 \label{p12-83-c20-s6:eq:heisenberg-9-nuevo}
\end{equation}
The commutator of two horizontal shifts is
\[
 [(a,b,0),(c,d,0)]=(0,0,bc-ad),
\]
which records \(-\sigma(u,v)\). Reducing each coordinate modulo three
gives
\begin{equation}
 H_3=\{(a,b,z):a,b,z\in\mathbb F_3\},\qquad |H_3|=27,
 \label{p12-83-c20-s6:eq:heisenberg-3-nuevo}
\end{equation}
with the same law. The nucleo of
\(\mathcal H_9\to H_3\) is
\((3\mathbb Z/9\mathbb Z)^3\), of cardinal 27.

\subsection{Weyl--Heisenberg representation and central phase}

The nonadic extension admits an explicit unitary realization. Let
\[
 \zeta_9:=\exp(2\pi i/9),
 \qquad
 \mathscr H_9^{\rm Sch}:=\mathbb C^9,
\]
based on \(\{|n\rangle:n\in\mathbb Z/9\mathbb Z\}\). We define
\begin{equation}
 X|n\rangle:=|n+1\rangle,
 \qquad
 Z|n\rangle:=\zeta_9^n|n\rangle.
 \label{p12-83-c20-s6:eq:weyl-desplazamiento-fase-u029}
\end{equation}
We have
\[
 ZX=\zeta_9XZ
\]
and, by multiplication,
\begin{equation}
 (X^aZ^b)(X^cZ^d)
 =\zeta_9^{bc}X^{a+c}Z^{b+d}.
 \label{p12-83-c20-s6:eq:cociclo-weyl-nonadico-u029}
\end{equation}

\begin{theorem}[Faithful Weyl--Heisenberg Representation]
\label{p12-83-c20-s6:thm:representacion-weyl-heisenberg-u029}
The map
\begin{equation}
 \rho_{\rm WH}:\mathcal H_9\longrightarrow
 U(\mathscr H_9^{\rm Sch}),
 \qquad
 \rho_{\rm WH}(a,b,z):=\zeta_9^zX^aZ^b,
 \label{p12-83-c20-s6:eq:representacion-weyl-heisenberg-u029}
\end{equation}
is a faithful unitary representation. For
\(u=(a,b)\) and \(v=(c,d)\), its commutator satisfies
\begin{equation}
 [\rho_{\rm WH}(u,0),\rho_{\rm WH}(v,0)]
 =\zeta_9^{\,bc-ad}I
 =\zeta_9^{-\sigma(u,v)}I.
 \label{p12-83-c20-s6:eq:conmutador-weyl-holonomia-u029}
\end{equation}
\end{theorem}

\begin{proof}
The identity \eqref{p12-83-c20-s6:eq:cociclo-weyl-nonadico-u029} reproduces the term
\(bc\) of \eqref{p12-83-c20-s6:eq:heisenberg-9-nuevo}; hence the
homomorphism property. If \(\zeta_9^zX^aZ^b=I\), the action on the basis successively forces
\(a=0\), \(b=0\), and \(z=0\). This proves faithfulness. Comparing
the products in the two orders gives
\(\zeta_9^{bc-ad}\), which is the alternating phase of
\eqref{p12-83-c20-s6:eq:forma-alternante-e6-nueva}.
\end{proof}

The representation distinguishes two conjugate observables. Let
\[
 Q_{\rm H}:=\operatorname{diag}(0,1,\ldots,8),
 \qquad
 F_{jk}:=\frac13\zeta_9^{jk},
 \qquad
 P_{\rm H}:=F^\dagger Q_{\rm H}F.
\]
The eigenbases of \(Q_{\rm H}\) and \(P_{\rm H}\) are mutually
unbiased, since \(|F_{jk}|^2=1/9\). For every unit vector
\(\psi\in\mathscr H_9^{\rm Sch}\),
\begin{align}
 \operatorname{Var}_\psi(Q_{\rm H})
 \operatorname{Var}_\psi(P_{\rm H})
 &\ge
 \frac14\left|
 \langle\psi,[Q_{\rm H},P_{\rm H}]\psi\rangle
 \right|^2,
 \label{p12-83-c20-s6:eq:incertidumbre-weyl-u029}\\
 H_{Q_{\rm H}}(\psi)+H_{P_{\rm H}}(\psi)
 &\ge\log9.
 \label{p12-83-c20-s6:eq:incertidumbre-entropica-weyl-u029}
\end{align}
These operators are dimensionless. A physical unit chart is a
later realization and does not intervene in the central extension.

\subsection{Action of \(10\) displazamientos over the grafo of Cayley}

For a linear functional
\(\ell(a,b)=pa+qb\), set
\[
 D_\ell(a,b)=\left(a,b,\frac{ab}{2}+\ell(a,b)\right),
 \qquad Z=(0,0,1),
\]
where \(2^{-1}=2\) in \(\mathbb F_3\). The inverse-closed set
\begin{equation}
 S_\ell=
 \{D_\ell(v):v\in\mathbb F_3^2\setminus\{0\}\}
 \cup\{Z,Z^{-1}\}
 \label{p12-83-c20-s6:eq:conjunto-conexion-e6}
\end{equation}
has ten elements and does not contain the identity.

\begin{theorem}[Independence of the linear scale]
\label{p12-83-c20-s6:thm:independencia-calibre-e6}
The nine functionals \(\ell\) produce isomorphic Cayley graphs. The
nine sets \(S_\ell\) form a single orbit under
\(\operatorname{Aut}(H_3)\).
\end{theorem}

\begin{proof}
The map
\[
 \phi_\ell(a,b,z)=(a,b,z+\ell(a,b))
\]
is a central automorphism and satisfies
\(\phi_\ell(S_0)=S_\ell\). To write the complete family in the
polarized coordinates of \cref{p12-83-c20-s6:eq:heisenberg-9-nuevo}, let
\(c(v,w)=b c\) for \(v=(a,b)\), \(w=(c,d)\). Given
\(M\in\operatorname{GL}(2,3)\), the defect
\[
 B_M(v,w)=c(Mv,Mw)-\det(M)c(v,w)
\]
is symmetric, because the alternating parts of the two terms coincide. We define the quadratic correction
\begin{equation}
 q_M(v)=2B_M(v,v),
 \qquad
 q_M(v+w)-q_M(v)-q_M(w)=B_M(v,w),
 \label{p12-83-c20-s6:eq:correccion-cuadratica-automorfismos-h3}
\end{equation}
where \(2=2^{-1}\) in \(\mathbb F_3\). The complete family of
automorphisms is then written
\begin{equation}
 (v,t)\longmapsto
 \bigl(Mv,\det(M)t+q_M(v)+\ell(v)\bigr),
 \qquad M\in\operatorname{GL}(2,3),\
 \ell\in(\mathbb F_3^2)^*.
 \label{p12-83-c20-s6:eq:automorfismos-h3-completos}
\end{equation}
The polar identity of \cref{p12-83-c20-s6:eq:correccion-cuadratica-automorfismos-h3}
directly proves that the central product is preserved. The family has
\(48\cdot9=432\) elements. Exhaustive enumeration of the
\(\binom{13}{5}=1287\) inverse-closed subsets of cardinality ten
finds exactly those nine with the parameters derived
below.
\end{proof}

Write \(\underline S=\sum_{s\in S}s\) in
\(\mathbb Z[H_3]\). Counting ordered products gives
\begin{equation}
 \underline S^{\,2}
 =10e+\underline S+5(\underline{H_3}-e-\underline S)
 =5\underline{H_3}+5e-4\underline S.
 \label{p12-83-c20-s6:eq:identidad-grupo-schlafli-nueva}
\end{equation}
Indeed, the identity receives ten factorizations; an element of
\(S\) receives one; each of the other sixteen receives five.

\begin{theorem}[Graph of the twenty-seven lines]
\label{p12-83-c20-s6:thm:grafo-27-rectas-nuevo}
The matrix of adjacency \(A_\ell\) of
\(\Gamma_\ell=\operatorname{Cay}(H_3,S_\ell)\) satisfies
\begin{equation}
 A_\ell^2=5I-4A_\ell+5J.
 \label{p12-83-c20-s6:eq:identidad-schlafli-nueva}
\end{equation}
Therefore,
\[
 \Gamma_\ell=\operatorname{srg}(27,10,1,5),
 \qquad
 \operatorname{Spec}(A_\ell)=\{10^1,1^{20},(-5)^6\}.
\]
It is isomorphic to the intersection graph of the twenty-seven lines of a
smooth cubic surface.
\end{theorem}

\begin{proof}
The regular representation of
\cref{p12-83-c20-s6:eq:identidad-grupo-schlafli-nueva} produces
\cref{p12-83-c20-s6:eq:identidad-schlafli-nueva}. Its entries state that every vertex
has ten neighbors, that an adjacent pair has one common neighbor, and a
nonadjacent pair has five. On the constant line, \(A_\ell\) acts by 10;
on its complement, \(J=0\) and
\((A_\ell-I)(A_\ell+5I)=0\). Zero trace fixes the multiplicities 20 and 6.

In the model of a cubic obtained by blowing up six points of
\(\mathbb P^2\), the 27 classes are
\[
 E_i,\qquad L_{ij}=H-E_i-E_j,\qquad
 Q_i=2H-\sum_{j\ne i}E_j.
\]
With \(H^2=1\), \(E_i^2=-1\), \(H\cdot E_i=0\), and
\(E_i\cdot E_j=0\) for \(i\ne j\), its intersection matrix has the
same parameters. Identification with the classical configuration uses the
uniqueness theorem for this realization of the 27 lines
\cite{Manivel2005}; it is not attributed to a nonexistent certificate.
\end{proof}

The isomorphism assertion can be made fully effective by means of
a computed bijection. The labeling is not absolute: another linear section
or an automorphism of the surface produces another table in the same class.
\begin{longtable}{@{}c@{\,}c@{\,}c@{\qquad}l@{}}
\caption{Explicit bijection between \(H_3\) and the twenty-seven classes of
lines.}
\label{p12-83-c20-s6:tab:h3-27-rectas-u029}\\
\toprule
\(a\)&\(b\)&\(z\)&class of line\\
\midrule
\endfirsthead
\multicolumn{4}{@{}l}{\small\itshape Table \thetable\ (continued)}\\
\toprule
\(a\)&\(b\)&\(z\)&class of line\\
\midrule
\endhead
0&0&0&\(E_1\)\\
0&0&1&\(L_{12}\)\\
0&0&2&\(Q_2\)\\
0&1&0&\(L_{13}\)\\
0&1&1&\(Q_1\)\\
0&1&2&\(E_3\)\\
0&2&0&\(Q_3\)\\
0&2&1&\(E_2\)\\
0&2&2&\(L_{23}\)\\
1&0&0&\(L_{14}\)\\
1&0&1&\(L_{35}\)\\
1&0&2&\(L_{26}\)\\
1&1&0&\(L_{36}\)\\
1&1&1&\(L_{24}\)\\
1&1&2&\(L_{15}\)\\
1&2&0&\(L_{25}\)\\
1&2&1&\(L_{16}\)\\
1&2&2&\(L_{34}\)\\
2&0&0&\(Q_4\)\\
2&0&1&\(L_{46}\)\\
2&0&2&\(E_6\)\\
2&1&0&\(E_5\)\\
2&1&1&\(Q_6\)\\
2&1&2&\(L_{56}\)\\
2&2&0&\(L_{45}\)\\
2&2&1&\(E_4\)\\
2&2&2&\(Q_5\)\\
\bottomrule
\end{longtable}

As a local control, the ten neighbors of \(e=(0,0,0)\) are the elements of
\(S_0\). The table transforms them into the ten classes that intersect \(E_1\):
the five \(L_{1j}\) and the five \(Q_j\), with \(2\le j\le6\). The
complete verification compares all 729 entries of both adjacency
matrices.

The complement \(B_\ell=J-I-A_\ell\) satisfies
\begin{equation}
 B_\ell^2=8I+2B_\ell+8J,
 \qquad
 \operatorname{srg}(27,16,10,8),
 \qquad
 \operatorname{Spec}(B_\ell)=\{16^1,4^6,(-2)^{20}\}.
 \label{p12-83-c20-s6:eq:complemento-schlafli}
\end{equation}
This is the Schläfli graph in the convention of complementary adjacency;
\(A_\ell\) is the intersection graph with parameters
\((27,10,1,5)\).

Each edge belongs to a unique triangle, since \(\lambda=1\). The number
of triangles is
\[
 \frac{27\cdot10}{6}=45.
\]
The full automorphism group has order
\[
 27\cdot1920=51840
\]
and, through its action on the line configuration, is identified with
\(W(E_6)\), not merely by equality of orders \cite{Manivel2005}.

\subsection{The root lattice realizing \texorpdfstring{\(E_6\)}{E6}}

To make explicit the object on which that group acts, let
\begin{equation}
 \operatorname{Pic}(X)
 =\mathbb ZH\oplus\bigoplus_{i=1}^6\mathbb ZE_i,
 \qquad
 H^2=1,\qquad E_i^2=-1,\qquad H\cdot E_i=E_i\cdot E_j=0\ (i\ne j),
 \label{p12-83-c20-s6:eq:picard-cubica-e6}
\end{equation}
the Picard lattice of the blow-up of six general points of
\(\mathbb P^2\). Its canonical class is
\[
 K_X=-3H+E_1+\cdots+E_6.
\]
The orthogonal complement
\begin{equation}
 L_{E_6}=K_X^\perp\subset\operatorname{Pic}(X),
 \qquad (x,y)_{E_6}:=-x\cdot y,
 \label{p12-83-c20-s6:eq:reticulo-raices-e6}
\end{equation}
is positive definite of rank six. A simple root basis is
\begin{equation}
 \begin{aligned}
 \alpha_1&=E_1-E_2,& \alpha_2&=E_2-E_3,&
 \alpha_3&=E_3-E_4,\\
 \alpha_4&=E_4-E_5,& \alpha_5&=E_5-E_6,&
 \alpha_6&=H-E_1-E_2-E_3.
 \end{aligned}
 \label{p12-83-c20-s6:eq:raices-simples-e6}
\end{equation}
Each \(\alpha_i\) is orthogonal to \(K_X\), has norm two, and its products
are \(-1\) exactly for adjacent pairs of the diagram
\(E_6\): the chain \(1-2-3-4-5\), with node \(6\) joined to node \(3\).
Therefore, its Gram matrix is the Cartan matrix of \(E_6\).

\begin{proposition}[The seventy-two roots and their reflections]
\label{p12-83-c20-s6:prop:setenta-dos-raices-e6}
The roots of \(L_{E_6}\) are
\begin{equation}
 \begin{aligned}
 &E_i-E_j &&(i\ne j),\\
 &\pm(H-E_i-E_j-E_k) &&(1\le i<j<k\le6),\\
 &\pm(2H-E_1-\cdots-E_6),
 \end{aligned}
 \label{p12-83-c20-s6:eq:enumeracion-raices-e6}
\end{equation}
in total \(30+40+2=72\). For each root \(\alpha\),
\begin{equation}
 s_\alpha(x)=x+(x\cdot\alpha)\alpha
 \label{p12-83-c20-s6:eq:reflexion-raiz-e6}
\end{equation}
preserves the intersection and \(K_X\). The six simple reflections generate
\(W(E_6)\) and act faithfully by permutations on the 27 classes
\(E_i,L_{ij},Q_i\) in \cref{p12-83-c20-s6:thm:grafo-27-rectas-nuevo}.
\end{proposition}

\begin{proof}
Substituting each class of \cref{p12-83-c20-s6:eq:enumeracion-raices-e6} into
\cref{p12-83-c20-s6:eq:picard-cubica-e6} gives \(K_X\cdot\alpha=0\) and
\(\alpha^2=-2\). The class count is as stated. The reflection formula
is the orthogonal reflection for the form \(-(\cdot)\); it preserves \(K_X\) because
\(K_X\cdot\alpha=0\). The simple system of
\cref{p12-83-c20-s6:eq:raices-simples-e6} generates all roots, and its Coxeter group is
\(W(E_6)\). Its action on the 27 classes of exceptional curves preserves
intersection, and therefore coincides with the action on the already
constructed graph \cite{Manivel2005}.
\end{proof}

\subsection{Stratification of the 729 Words}

We write a six-trit word as
\[
 w=(\ell_1,h_1,\ell_2,h_2,\ell_3,h_3),
\]
and we define two elements of \(H_3\),
\[
 L(w)=(\ell_1,\ell_2,\ell_3),\qquad
 H(w)=(h_1,h_2,h_3),
\]
using the central law in their products. According to
\(L(w)^{-1}H(w)\), the 729 words are divided into
\begin{align*}
 \mathcal D_\ell&=\{w:L(w)^{-1}H(w)=e\},\\
 \mathcal I_\ell&=\{w:L(w)^{-1}H(w)\in S_\ell\},\\
 \mathcal S_\ell&=\{w:L(w)^{-1}H(w)\notin\{e\}\cup S_\ell\}.
\end{align*}

\begin{theorem}[Hensel--Schläfli Partition]
\label{p12-83-c20-s6:thm:particion-hensel-schlafli-nueva}
\begin{equation}
 729=27+270+432,
 \label{p12-83-c20-s6:eq:particion-729-e6-nueva}
\end{equation}
corresponding to diagonal, incidence, and warping.
\end{theorem}

\begin{proof}
For each one of the 27 values of \(L(w)\), there is a unique value of
\(H(w)\) with identity difference, ten with difference in \(S_\ell\) and
sixteen remaining ones. \(27\) is multiplied by \(1,10,16\).
\end{proof}

The cardinal \(432\) of the warped stratum numerically coincides with the
multiplicity of the transverse region of \(\pi\) in
\cref{eq:descomposiciones-cinco-regiones-pi}. The two objects are not thereby identified: here
one counts words classified by \(L(w)^{-1}H(w)\), whereas there one
counts preimages of a register \(U_6\). Without an explicit and
equivariant intertwiner between those domains, \(432=432\) is only a cardinal control.

\subsection{Necessity of the central phase}

Enumerating all inverse-closed subsets of cardinality ten in the five groups of order 27 gives
\[
\begin{array}{c|c|c}
\text{group}&\text{abelian}&
\#\operatorname{srg}(27,10,1,5)\\ \hline
C_{27}&\text{yes}&0\\
C_9\times C_3&\text{yes}&0\\
C_3^3&\text{yes}&0\\
C_9\rtimes C_3&\text{no}&9\\
H_3&\text{no}&9.
\end{array}
\]
Thus, retaining 27 labels while abelianizing composition destroys the
configuration. Neither cardinality nor degree suffices.

If \(\mathcal T\) is the conjunto of 45 triangulos, the cubico
\[
 \mathcal C_\Gamma(x_1,\ldots,x_{27})
 =\sum_{\{i,j,k\}\in\mathcal T}x_ix_jx_k
\]
recovers adjacency, because every edge lies in a unique triangle. Its
permutation stabilizer is exactly
\[
 \operatorname{Aut}(\Gamma_\ell)\cong W(E_6).
\]

\begin{figure}[htbp]
\centering
\[
 \text{APP curvatures }\pm1
 \longrightarrow\sigma
 \longrightarrow\mathcal H_9
 \longrightarrow H_3
 \longrightarrow\Gamma_{27}
 \longrightarrow\mathcal T_{45}
 \longrightarrow W(E_6).
\]
\caption{Causal chain of the mixed-holonomy branch. It is parallel to
\(C_W\to W_{12}\to M_{12}\) and not a link of it.}
\label{p12-83-c20-s6:fig:cadena-holonomia-e6-u029}
\end{figure}

\begin{criterion}[Falsification criteria for branch \(E_6\)]
\label{p12-83-c20-s6:crit:refutacion-rama-e6-u029}
The branch is refuted if the curvatures are not opposite, if the central cocycle is erased, if the quadratic correction \eqref{p12-83-c20-s6:eq:correccion-cuadratica-automorfismos-h3} does not preserve the product, if \eqref{p12-83-c20-s6:eq:identidad-grupo-schlafli-nueva} does not produce the coefficients \((10,1,5)\), if an abelian group produces the same configuration, if the partition does not sum to 729, if the table \ref{p12-83-c20-s6:tab:h3-27-rectas-u029} does not intertwine the adjacency matrices, or if \(W(E_6)\) is identified only by its order and not by its action on the twenty-seven classes and the 72 roots.
\end{criterion}


\section{Oriented APP–Fock coupling and transverse rigidity of \(12\) fibres}
\label{p12-83-c20-s7:chap:acoplamiento-app-fock}
\index{APP--Fock!oriented coupling}
\index{Fock!degree two}
\index{sheet orientation}

The ordinary cyclotomic determinant does not distinguish \(C\) from \(C^{-1}\).
Sum--product information is preserved by explicitly introducing
sheet parity and an orientation line. The resulting coupling
produces index \(54\) from the APP aggregate. Once the four
coupling conditions have been fixed, \(54\), \(J\), and \(196884\) are not inputs to
the evaluation. This evaluative independence asserts neither universal uniqueness
of the ansatz nor historical independence of its design.

\subsection{Triangular mark on each orbit}

Let
\[
M=\mathbb Z[\mathbb Z/9\mathbb Z]
\]
and let \(M_{\mathrm{bal}}^{C_3}\) be the submodule of invariant elements
\(\nu\) whose six unit coefficients coincide:
\[
\nu_u=c(\nu)
\qquad
\bigl(u\in(\mathbb Z/9\mathbb Z)^\times\bigr).
\]
The coefficients of \(e_0,e_3,e_6\) permanecen libres.

For a free orbit
\(\mathcal O=\{u,4u,7u\}\), set
\[
b_v=e_v-e_{4v}\in A(\mathcal O).
\]
If \(b_v^\flat\) is the covector defined by the form of \(A_2\), the three
roots of the triangular frame satisfy
\[
\boxed{
\sum_{v\in\mathcal O}b_v\otimes b_v^\flat
=3I_{A(\mathcal O)}.}
\]
Consequently,
\[
\sum_{v\in\mathcal O}\nu_v\,b_v\otimes b_v^\flat
=3c(\nu)I.
\]
The factor three proceeds from the resolution of the identity by the root framework.

\subsection{Oriented coupling \(\mathcal H_{\Sigma,\Pi}^{(m)}\) between the balance \(C_3\), leaf parity, and degree \(2\) of Fock space}

For \(m\in2\mathbb N_{>0}\), let
\[
\begin{aligned}
\mathcal V_m&=\mathcal V_{\Sigma,m}\oplus\mathcal V_{\Pi,m},\\
\mathcal V_{\Sigma,m}&=(A_2\otimes\mathbb C)^{\oplus m/2},\qquad
\mathcal V_{\Pi,m}=(A_2\otimes\mathbb C)^{\oplus m/2},\\
g_m&=C^{\oplus m/2}\oplus(C^{-1})^{\oplus m/2},\\
\mathsf E_m&=I_{\mathcal V_{\Sigma,m}}
\oplus(-I_{\mathcal V_{\Pi,m}}).
\end{aligned}
\]
The specialization \(m=12\), with the labels of pair, sheet and orbit,
is \(\mathcal V_{12}=W_{\mathrm{APP}}\otimes\mathbb C\). This notation
prevents confusing the representational space with the Witt system
\(W_{12}=S(5,6,12)\). Degree two of the Fock space and the induced action are
\[
\mathcal F_2(\mathcal V_m)
=\mathcal V_m[2]\oplus\operatorname{Sym}^2(\mathcal V_m[1]),
\]
\[
\Gamma_2(g_m)
=g_m|_{\mathcal V_m[2]}
\oplus\operatorname{Sym}^2(g_m|_{\mathcal V_m[1]}).
\]

Let \(o_{\Sigma\Pi}\) be the signed orientation line
\(\Pi-\Sigma\), and let \(M_{\mathrm{bal}}^{C_3,-}\) be the oriented copy in
which the exchange of leaves acts by \(\nu\mapsto-\nu\). We define
\[
\boxed{
\begin{aligned}
\mathcal H_{\Sigma,\Pi}^{(m)}&:
M_{\mathrm{bal}}^{C_3,-}\longrightarrow
\operatorname{End}(\mathcal F_2(\mathcal V_m))\otimes o_{\Sigma\Pi},\\
\mathcal H_{\Sigma,\Pi}^{(m)}(\nu)
&=3c(\nu)
\left(0_{\mathcal V_m[2]}\oplus\operatorname{Sym}^2\mathsf E_m\right)
\otimes o_{\Sigma\Pi}.
\end{aligned}}
\]
The transformation is linear in \(c\), is supported on
\(\operatorname{Sym}^2(\mathcal V_m[1])\), uses leaf parity, and adopts the
normalization of the triangular frame.

Let \(\sigma\) the exchange of sheets. Then
\[
\nu\longmapsto-\nu,\qquad
o_{\Sigma\Pi}\longmapsto-o_{\Sigma\Pi},\qquad
\mathsf E_m\longmapsto-\mathsf E_m,
\]
and
\[
\operatorname{Sym}^2(-\mathsf E_m)
=\operatorname{Sym}^2(\mathsf E_m).
\]
Therefore,
\[
\mathcal H_{\Sigma,\Pi}^{(m)}(\sigma\nu)
=\sigma\mathcal H_{\Sigma,\Pi}^{(m)}(\nu).
\]
When the orientation is reversed, the dual covector also changes sign; thus
\((-o_{\Sigma\Pi}^{\vee})(-o_{\Sigma\Pi})=1\), and the oriented evaluation
remains normalized.

Fixing
\(o_{\Sigma\Pi}^{\vee}(o_{\Sigma\Pi})=1\), the oriented trace is
\[
\operatorname{Tr}_{o,\mathcal F_2}
(T\otimes o_{\Sigma\Pi})
=\operatorname{Tr}_{\mathcal F_2}(T).
\]

\begin{theorem}[APP-oriented Index--Fock]
\label{p12-83-c20-s7:thm:indice-orientado-app-fock}
Suppose that \(m\) fibers \(A_2\), with \(m\) even, are distributed equally
between the orientations \(C\) and \(C^{-1}\). Then
\[
\boxed{
\operatorname{Tr}_{o,\mathcal F_2}
\left(
\mathcal H_{\Sigma,\Pi}^{(m)}(\nu)\Gamma_2(g_m)
\right)
=-\frac{3m\,c(\nu)}2.}
\]
For the aggregate produced by APP,
\[
\delta=\nu_\Pi-\nu_\Sigma
=12e_0+3e_3+3e_6
-3\sum_{u\in(\mathbb Z/9\mathbb Z)^\times}e_u,
\]
one has \(c(\delta)=-3\) and
\[
m=3\ \text{pairs}\cdot2\ \text{sheets}\cdot2\ \text{orbits}=12.
\]
Therefore,
\[
\boxed{
\operatorname{Tr}_{o,\mathcal F_2}
\left(
\mathcal H_{\Sigma,\Pi}^{(12)}(\delta)\Gamma_2(g_{12})
\right)=54.}
\]
\end{theorem}

\begin{proof}
Let \(A=\mathsf E_mg_m\). The contributions of the two leaves cancel:
\(\operatorname{tr}A=0\). Since \(\mathsf E_m\) commutes with \(g_m\) and
\(\mathsf E_m^2=I\),
\[
\operatorname{tr}(A^2)=\operatorname{tr}(g_m^2)=-m.
\]
The identity
\[
\operatorname{tr}(\operatorname{Sym}^2A)
=\frac{(\operatorname{tr}A)^2+\operatorname{tr}(A^2)}2
\]
gives \(-m/2\); the frame supplies \(3c(\nu)\).
\end{proof}

\begin{corollary}[Minimality of the Bosonic Degree Detector]
\label{p12-83-c20-s7:cor:minimalidad-grado-dos-app-fock-u030}
In the declared class of couplings, degree one does not detect the oriented
aggregate, and degree two is the first positive bosonic degree that
detects it:
\[
 \operatorname{tr}(\mathsf E_mg_m)=0,
 \qquad
 \operatorname{tr}\!\left(
 \operatorname{Sym}^2(\mathsf E_mg_m)
 \right)=-\frac m2\neq0.
\]
\end{corollary}

\begin{proof}
The first equality is the cancellation of the two sheets. The second was
obtained in the proof of \cref{p12-83-c20-s7:thm:indice-orientado-app-fock}. Since degree
zero contains only the unit and degree one has zero trace, degree two is
the first with a nonzero response.
\end{proof}

\subsection{Transverse rigidity of the number of fibers}

For an integer \(m\ge1\), we define the formal series
\begin{equation}
 T_m(q):=
 q^{-1}\prod_{n\ge1}(1+q^n+q^{2n})^{-m}
 \in q^{-1}\mathbb Z[[q]]
 \label{p12-83-c20-s7:eq:familia-fock-m-fibras-u030}
\end{equation}
and, in the orthogonal sum \(A_2^m\), the radial vector
\begin{equation}
 v_m:=4\rho_1+\rho_2+\cdots+\rho_m,
 \qquad
 \|\rho_j\|^2=2.
 \label{p12-83-c20-s7:eq:familia-radial-m-fibras-u030}
\end{equation}
We construct two entire functions without using \(m=12\) beforehand:
\[
 F(m):=[q]T_m(q),
 \qquad
 R(m):=\|v_m\|^2.
\]

\begin{proposition}[Fock coefficient and radial norm]
\label{p12-83-c20-s7:prop:dos-funciones-rigidez-doce-u030}
For every \(m\ge1\),
\begin{equation}
 F(m)=\frac{m(m-3)}2,
 \qquad
 R(m)=2(m+15).
 \label{p12-83-c20-s7:eq:funciones-rigidez-doce-u030}
\end{equation}
\end{proposition}

\begin{proof}
Up to second order,
\[
 (1+q+q^2)^{-m}
 =1-mq+\frac{m(m-1)}2q^2+O(q^3),
\]
\[
 (1+q^2+q^4)^{-m}=1-mq^2+O(q^4).
\]
The factors with \(n\ge3\) do not contribute to the coefficient of \(q^2\).
Since the exterior factor of \(T_m\) is \(q^{-1}\),
\[
 F(m)=\frac{m(m-1)}2-m=\frac{m(m-3)}2.
\]
By orthogonality of the \(m\) copies of \(A_2\),
\[
 R(m)=4^2\|\rho_1\|^2+
 \sum_{j=2}^{m}\|\rho_j\|^2
 =32+2(m-1)=2(m+15).
\]
\end{proof}

\begin{theorem}[Cyclotomic--radial rigidity]
\label{p12-83-c20-s7:thm:rigidez-transversal-doce-u030}
The equality \(F(m)=R(m)\) is equivalent to
\begin{equation}
 (m-12)(m+5)=0.
 \label{p12-83-c20-s7:eq:ecuacion-rigidez-doce-u030}
\end{equation}
In the domain \(m\in\mathbb Z_{>0}\), the unique solution is
\[
 \boxed{m=12,\qquad F(12)=R(12)=54.}
\]
\end{theorem}

\begin{proof}
Equating the two expressions of
\eqref{p12-83-c20-s7:eq:funciones-rigidez-doce-u030} produce
\[
 m(m-3)=4(m+15),
\]
that is,
\[
 m^2-7m-60=(m-12)(m+5)=0.
\]
The second root is negative and lies outside the domain.
\end{proof}

This rigidity belongs to the declared cyclotomic--radial family. It does not use the
value \(54\) to choose \(m\): it first constructs \(F(m)\) and \(R(m)\), and
then resolves their equality.

\subsection{Factorization of the APP–Fock coupling: hypotheses, invariants, necessary
conditions for validity, and obstructions}

The coupling factorizes by
\[
M_{\mathrm{bal}}^{C_3,-}/\ker c
\cong\mathbb Z\otimes o_{\Sigma\Pi},
\]
and the induced map is injective because
\(\operatorname{Sym}^2\mathsf E_m\neq0\).
The formula is selected by the four declared conditions:
linearity in \(c\), support in
\(\operatorname{Sym}^2(\mathcal V_m[1])\), sheet parity,
and normalization by the triangular frame. No uniqueness is asserted
among all natural couplings, nor historical independence of the
ansatz. Once those conditions have been fixed, \(54\), \(J\), and \(196884\) are not
inputs to the evaluation.

Three controls show that the result is not determined by a retroactive normalization:
\[
(m,c)=(12,-2)\longmapsto36,\qquad
(m,c)=(10,-3)\longmapsto45,
\]
and
\[
\delta_k=\delta+ke_0
\]
leaves \(\mathcal H_{\Sigma,\Pi}^{(12)}\) invariant while the digital weight
changes from \(54\) to \(54+9k\). In particular, for \(k=1\),
\[
\operatorname{Tr}_{o,\mathcal F_2}
\left(
\mathcal H_{\Sigma,\Pi}^{(12)}(\delta+e_0)\Gamma_2(g_{12})
\right)=54,
\qquad
w(\delta+e_0)=63.
\]

\begin{proposition}[Need for orientation data]
\label{p12-83-c20-s7:prop:orientacion-app-fock}
If \(o_{\Sigma\Pi}\) is forgotten, every natural scalar that is linear in
\(\nu_\Pi-\nu_\Sigma\) and depends only on the class of the representation
vanishes.
\end{proposition}

\begin{proof}
\(C\) and \(C^{-1}\) are conjugate, so the representational class is
even under sheet exchange. The aggregate
\(\nu_\Pi-\nu_\Sigma\) is odd. Without the orientation line, a
transformation that is simultaneously even and odd is zero.
\end{proof}

The theorem thus establishes the exact transport of the aggregate to Fock
degree two. The local state-by-state defect preserves the descent obstruction
of \cref{p12-83-c1-s5:cor:obstruccion-isotipica-app-autonoma}; the two statements belong to
distinct domains and are compatible.


% Unidad U031. Alineamiento APP--Witt, automorfia de orden tres,
% construccion VOA y Moonshine.
